\documentclass[12pt]{article}

\usepackage[utf8]{inputenc}
\usepackage{CJKutf8}
\usepackage{amsmath,amssymb}
\usepackage{amsthm}
\usepackage{geometry}
\geometry{
  a4paper,
  top=25mm,
  bottom=25mm,
  left=20mm,
  right=20mm
}
\usepackage{xcolor}
\usepackage{url}
\usepackage[unicode,colorlinks=true,linkcolor=blue,urlcolor=blue]{hyperref}
\usepackage{xurl}
\Urlmuskip=0mu plus 2mu
\sloppy
\usepackage{tocloft}

% 目录中 section 条目使用点线填充到页码
\renewcommand{\cftsecleader}{\cftdotfill{\cftdotsep}}
\renewcommand{\refname}{参考文献}

% 基本定理环境（工作笔记：形式系统风格）
\newtheorem{definition}{定义}[section]
\newtheorem{axiom}{公理}[section]
\newtheorem{assumption}{假设}[section]
\newtheorem{theorem}{定理}[section]
\newtheorem{proposition}{命题}[section]
\newtheorem{corollary}{推论}[section]
\newtheorem{lemma}{引理}[section]
\newtheorem{remark}{备注}[section]
\newtheorem{Certificate}{证书}[section]

% 记号（仅用于本笔记自洽引用，避免依赖主稿宏定义）
\newcommand{\code}[1]{\path{#1}}
\newcommand{\GFramework}{\textsf{G-Framework}}

% 避免少量不可控的换行溢出（尤其是混排 CJK 与等宽字体时）

\hfuzz=700pt
\hbadness=10000

\begin{document}
\begin{CJK*}{UTF8}{gkai}

\title{语义函子：正则生成、裁决与规范化接口\\（工作笔记：形式系统版）}
\author{Gao Zheng（高政）}
\date{2026-01-14}
\maketitle

\section*{前言}
\addcontentsline{toc}{section}{前言}

本文的目标不是单独讨论语义函子、正则生成、裁决或规范化接口，而是在
\GFramework{} 的语义工程链条中，把表达输入、标准正则、MDQ 分割、外包证书、
粒度审计、多次扫描与固定点收敛组织成一个可复验的语义规范化系统。若这些对象
分散出现，语义函子容易被看成抽象映射，正则生成容易被看成字符串技巧，可信性也
容易停留在主观判断；本文将其改写为可裁决、可留痕、可回放的工程接口。

本文的第一层立场是语义函子的工程落地。语义函子在本文中不只是从表达到意义的
抽象箭头，而是从自然语言或半结构表达进入标准正则、候选同伦类与验证接口的
可调用组件。它把“表达是否合格”转化为“能否生成可裁决的标准表示并进入证书流”的
问题。

本文的第二层立场是可信性的降级。本文不把“可信”作为不可再分的整体标签，而是
把它分解为 MDQ 粒度、外包证书、复验日志、扫描历史与固定点判据。由此，可信性
不再依赖单次判断，而依赖可追溯的最小单元、可重复的验证路径与候选集合的单调
收缩。

本文的第三层立场是规范化的鲁棒收敛。多次扫描不是重复劳动，而是让候选同伦类
集合逐步收缩，直到达到固定点或暴露出不可判定边界。规范化因此不仅产生标准文本，
还产生一个“为什么当前标准足够稳定”的证书过程。

本文的第四层立场是 LLM 与确定性裁决的分工。LLM 可以作为候选生成器，但不能成为
最终正确性的口头担保。本文把生成与裁决分离：生成阶段负责提出候选正则或候选
解释，裁决阶段通过编译、样例验证、复杂度阈值、变形测试和回归日志决定其是否
进入标准集。

本文的第五层立场是 MDQ 的可移植层意义。MDQ 将表达拆成最小可审计单元，使外包
证书、复验、同伦类归属和语义迁移都能在细粒度上发生。由此，语义函子不再依赖
整体句子的不可分判断，而可以把复杂表达分解为可复验的小块，再通过组合恢复整体。

本文的第六层立场是多次扫描的单调收缩。候选集合如果每轮都可能任意扩大，就无法
形成稳定规范；本文要求扫描过程在证书规则下呈现单调收缩或受控更新，直到达到固定
点。固定点不是绝对真理，而是给定规则、样例和复杂度边界下的稳定规范化状态。

本文的第七层立场是为后续命中和治理提供语义入口。无论是开放系统治理、策略簇
生成还是 PDE 残差命中，都需要可靠输入表达。本文将这种输入可靠性前置为语义函子
和正则规范化问题，使后续数学或工程系统不必直接承受未经裁决的自然语言不确定性。

因此，本文在整体 \GFramework{} 中承担的是“语义函子到正则规范化接口”的作用。
它向前连接语义规范场与边缘算子，向中间连接 MDQ、外包证书与固定点扫描，向后
为开放系统治理、同伦命中和 PDE 残差范式提供语义输入的可裁决前端。概括地说，
本文把“表达是否可信”推进为“语义函子生成、正则规范化、证书复验与鲁棒固定点”的
形式系统。

更具体地说，本文把“自然语言表达”降级为需要通过接口的输入对象。输入可以由模型
生成，也可以由人给出，但进入系统之前必须经过规范化、验证和鲁棒性测试。这样，
后续形式系统不再直接承受表达歧义，而只接收已裁决的标准正则或明确失败符号。

从方法论上看，本文给出了语义工程中的“确定性层”。不确定生成负责提出候选，确定性
检测负责裁决候选，扫描收敛负责稳定候选集合，证书日志负责追溯每一次接受或拒绝。
这正是后续统计命中中“随机生成--确定性判定”结构在语义输入层的提前出现。

从整体谱系看，本文把 \GFramework{} 的语义入口工程化。没有这个入口，后续开放系统
治理、策略簇构造和 PDE 残差命中都可能被含混表达污染。本文的意义在于先把表达
收束为可裁决对象，再让形式系统在此基础上运行。

\setcounter{tocdepth}{3}
\setcounter{secnumdepth}{3}
\tableofcontents
\clearpage

%%%%%%%%%%%%%%%%%%%%%%%%%%%%%%%%%%%%%%%%%%%%%%%%%%%%%%%%%%%%
\section{把“语义函子”落到工程：从表达输入到可裁决的标准正则}
\label{sec:semantic-functor-regex-interface}
%%%%%%%%%%%%%%%%%%%%%%%%%%%%%%%%%%%%%%%%%%%%%%%%%%%%%%%%%%%%

\begin{remark}[核心陈述（工程化接口）]
把您说的“语义函子”落到工程里，本质上就是做一个可复用映射器：
\begin{itemize}
  \item 输入：一段表达（通常是自然语言描述、字段含义描述、或一类输入样例）；
  \item 输出：要么失败符号 $\bot$（表示：在您定义的规则边界内，无法给出安全、确定的正则），
  要么返回一个已归一化的标准正则 $r$（表示：可给出、且已规范化为您认可的“标准形式”）。
\end{itemize}
关键在于：LLM 负责“生成候选”，但“是否可接受、是否鲁棒”必须由确定性的校验/测试框架裁决，
而不是继续让 LLM 口头保证。
这体现了一种“专注于可控”的工程取向：把不确定的生成收敛为可复验的裁决与回归闭环，而非把正确性外包给模型叙述\cite{gaozheng-analytic-ai}。
\end{remark}

下面给出一套可直接落地的构造方案，并在第\ref{subsec:examples-true-false-boundary}小节给出“\code{ok:true}+标准正则 / \code{ok:false}”的边界示例。

\begin{remark}[备注：LLM API 快速构造与鲁棒性检测闭环]
在实现层，可将 LLM 的 API 用作候选生成器，从而快速构造语义函子接口：输入表达 $x$，
输出要么为失败（\code{ok:false}，对应失败符号 $\bot$），要么为通过并返回规范化后的标准正则（\code{ok:true} 且携带 \code{regex}）。
但对外接口的“可信”必须由确定性的检测机制给出：通过白名单/编译/样例验证/复杂度阈值与变形测试回归（第\ref{subsec:metamorphic-testing}小节），
检验该语义函子在语义等价改写与噪声扰动下的稳定性与可控边界。
\end{remark}

\subsection{合同：把通过域/失败域写死}
\label{subsec:contract}

\begin{definition}[语义函子接口（判定 + 标准化正则）]
令 $\mathsf{Expr}$ 为输入表达域。令 $\mathsf{Re}_{\mathrm{safe}}$ 为目标正则语言的“安全子集”，并令
$\mathsf{Re}_{\mathrm{can}}\subseteq \mathsf{Re}_{\mathrm{safe}}$ 为其规范形集合（canonical forms）。
称映射
\[
F:\mathsf{Expr}\to \mathsf{Re}_{\mathrm{can}}\sqcup\{\bot\}
\]
为本章的“语义函子接口”，并约定：
\begin{enumerate}
  \item $F(x)=\bot$ 当且仅当：输入表达 $x$ 的语义不确定、范围过大、或需要超出安全正则子集的表达能力
  （例如需要上下文相关约束、跨字段依赖、或需要 \emph{lookbehind/backreference} 等高风险特性）；
  \item $F(x)=r$（其中 $r\in\mathsf{Re}_{\mathrm{can}}$）当且仅当：能给出一个属于 $\mathsf{Re}_{\mathrm{safe}}$ 的正则，
  且该正则已通过规范化（canonicalization）被归一到您认可的标准形式。
\end{enumerate}
\end{definition}

\begin{remark}[强烈建议：限定“安全正则子集”】【工程可控性】]
为了让“鲁棒性测试”与运行时资源消耗可控，建议把 $\mathsf{Re}_{\mathrm{safe}}$ 明确限定为 \textsf{RE2} 风格（或等价子集）\cite{cox-re2,cox-regex}，并禁止下列高风险特性：
\begin{itemize}
  \item backreference（如 \verb|\1| 等）；
  \item 任意 lookaround（尤其是负向 lookbehind）；
  \item 可能导致灾难性回溯的结构（例如嵌套可变长量词）。
\end{itemize}
在该约束下，“标准正则”对应真正意义上的正则语言，测试与运行更稳定。
\end{remark}

\subsection{输出格式：固定为机器可检验的 JSON 合同}
\label{subsec:json-contract}

外部接口建议采用严格 JSON\cite{rfc8259} 作为序列化格式，并且只承认 \code{ok} 与 \code{regex} 两个字段（其余字段仅作内部调试与回归辅助）。

\paragraph{成功：}
\begin{verbatim}
{"ok": true, "regex": "^(?:...)$"}
\end{verbatim}

\paragraph{失败：}
\begin{verbatim}
{"ok": false}
\end{verbatim}

\begin{remark}[允许扩展字段（不影响外部合同）]
为便于“可调试”与“可做鲁棒性证书”，内部可以允许扩展字段（外部只认 \code{ok} 与 \code{regex}）：
\begin{verbatim}
{
  "ok": true,
  "regex": "^(?:...)$",
  "positive_examples": ["...","..."],
  "negative_examples": ["...","..."],
  "notes": "简短说明：哪些边界被排除"
}
\end{verbatim}
\end{remark}

\subsection{两段式流水线：LLM 生成候选 + 确定性校验器裁决}
\label{subsec:two-stage-pipeline}

\begin{definition}[两段式组合接口]
将系统拆分为：
\begin{enumerate}
  \item 候选生成器 $G_{\mathrm{LLM}}$：给定 $x\in\mathsf{Expr}$，生成候选 JSON（可能包含候选正则与自证样例）；
  \item 确定性校验器 $V$：对候选进行白名单扫描、编译与样例验证，并在通过时做规范化。
\end{enumerate}
最终对外的语义函子接口可被实现为复合：
\[
F(x):=V(G_{\mathrm{LLM}}(x))\in \mathsf{Re}_{\mathrm{can}}\sqcup\{\bot\}.
\]
\end{definition}

\subsubsection{LLM 侧：生成候选（受控）}

让 LLM 严格做三件事（缺一不可）：
\begin{enumerate}
  \item 给出候选正则（限制在您规定的安全子集内）；
  \item 给出若干正例/反例（用于自证一致性与后续回归）；
  \item 在不确定时必须输出 \code{"ok": false}（而不是硬编）。
\end{enumerate}

为了减少漂移并提升可控性，实践上建议：
\begin{itemize}
  \item 温度设为 $0$ 或极低；
  \item 强制 JSON 输出（结构固定）；
  \item 在提示词中明确写出“失败条件”（什么时候必须返回失败）。
\end{itemize}

\subsubsection{确定性校验器：您真正的“裁决机制”}

校验器只做机器可判定的事情，例如：
\begin{itemize}
  \item JSON 可解析；
  \item \code{regex} 可编译；
  \item 正则仅使用允许的语法（白名单扫描）；
  \item 必须带 \verb|^| 与 \verb|$|（强制全匹配，防止“局部匹配误过”）；
  \item 对 \code{positive_examples} 必须全部匹配；
  \item 对 \code{negative_examples} 必须全部不匹配；
  \item 长度/复杂度阈值（避免过长正则造成维护困难或资源风险）。
\end{itemize}
任一检查失败：直接判 $F(x)=\bot$，并记录失败原因（用于迭代提示词或补充训练样例/回归集）。

\subsection{鲁棒性：用变形测试（metamorphic testing）做自动化回归}
\label{subsec:metamorphic-testing}

您要的“测试是否满足该语义函子的鲁棒性”，可用不依赖人工大规模标注的变形测试（metamorphic testing）\cite{chen-metamorphic}构造自动回归。
下面给出三类常用测试口径。

\subsubsection{语义等价输入的稳定性：同义改写不应改变输出}

对同一语义表达 $x$，生成一组等价变体 $\{x_i\}$（例如同义替换、中文/英文混写、标点/空白扰动、口语化/正式化改写等），
比较 $F(x_i)$ 的一致性。可用从严到宽的分档判据：
\begin{itemize}
  \item 强鲁棒：规范化后输出正则字符串完全一致；
  \item 中鲁棒：字符串不完全一致，但在抽样集上行为等价（见下小节）；
  \item 弱鲁棒：至少都返回通过，且都能通过自给例子验证（适合早期）。
\end{itemize}

\subsubsection{噪声不敏感性：输入样本扰动，输出应稳定}

若输入表达为“样本驱动”（例如给几条样例让系统归纳正则），则可对样本施加扰动：
大小写、前后空格、全角半角、常见 OCR 错误（\code{O/0}、\code{l/1}）等，
观察 $F$ 的输出是否稳定，或是否以可解释方式收敛到同一规范形。

\subsubsection{行为等价性：用抽样集近似检验两个 regex 的等价}

严格证明两个正则等价可诉诸自动机最小化\cite{hopcroft-ullman}；工程上通常用“抽样近似判据”足够实用。
对两个输出正则 $r_1,r_2$，构造测试字符串集合 $S$（随机串 + 边界串 + 生成串），检查：
\[
\forall s\in S,\quad \mathrm{match}(r_1,s)=\mathrm{match}(r_2,s).
\]
集合 $S$ 的构造建议分层：
\begin{itemize}
  \item 边界：最短、最长、临界字符（如 \code{-}、\code{_}、\code{.}、\code{@} 等）；
  \item 随机：按字符集与长度分布采样；
  \item 对抗：注入容易误判的串（例如在邮箱里插入两个 \code{@}，或在日期里插入 \code{13} 月）。
\end{itemize}
该方法可快速捕捉“提示词/模型/规则集变更”导致的语义漂移。

\subsection{规范化（canonicalization）：让“标准正则”真正标准}
\label{subsec:canonicalization}

\begin{definition}[规范化算子]
定义一个规范化算子
\[
\mathrm{can}:\mathsf{Re}_{\mathrm{safe}}\to \mathsf{Re}_{\mathrm{can}},
\]
把等价但写法不同的正则统一到可比较、可回归的标准形式。
\end{definition}

可选的规范化规则例如：
\begin{itemize}
  \item 强制全匹配：统一补齐 \verb|^...$|；
  \item 强制使用非捕获组：统一为 \verb|(?:...)|；
  \item 字符类排序：把 \verb|[A-Za-z0-9_]| 之类字符类排成固定顺序；
  \item 统一替代写法：例如 \verb|\d| 与 \verb|[0-9]| 统一选一种；
  \item 去掉冗余括号、冗余分支；
  \item 限定量词风格：统一用 \verb|{m,n}| 或统一用 \verb|+|/\verb|*|（二选一并固化）。
\end{itemize}
完成规范化后，“同义输入稳定性”更容易达到“强鲁棒”（字符串一致）。

\subsection{最小可落地路径：从白名单到 CI 的六步}
\label{subsec:minimal-implementation-path}

\begin{enumerate}
  \item 选定目标正则方言与白名单（推荐 \textsf{RE2} 子集）；
  \item 定义输出 JSON 合同：只对外承诺 \code{ok} + \code{regex}；
  \item 编写提示词：强制 JSON；明确失败条件；要求给正反例；
  \item 实现校验器：编译、白名单、例子验证、复杂度阈值；
  \item 实现鲁棒性测试器：生成等价变体 $\{x_i\}$、调用 $F(x_i)$、做规范化与一致性/行为等价抽样检验；
  \item 接入 CI：每次改提示词、换模型、或改规则集，都跑回归并输出“鲁棒性报告”（通过率、漂移样例、失败原因聚类）。
\end{enumerate}

\subsection{边界示例：对齐 “True+标准正则 / False”}
\label{subsec:examples-true-false-boundary}

\subsubsection{例 A：语义明确且可正则化（返回通过）}

\paragraph{输入表达：}
\begin{quote}
匹配标准 UUID（含连字符）
\end{quote}

\paragraph{输出（示例）：}
\begin{verbatim}
{"ok": true, "regex": "^[0-9a-fA-F]{8}-[0-9a-fA-F]{4}-[0-9a-fA-F]{4}-[0-9a-fA-F]{4}-[0-9a-fA-F]{12}$"}
\end{verbatim}

\subsubsection{例 B：语义超出正则能力或高度依赖上下文（返回失败）}

\paragraph{输入表达：}
\begin{quote}
匹配一个合法的 JSON（任意嵌套深度）；

匹配括号完全配对且层数无限的表达式。
\end{quote}

\paragraph{输出（示例）：}
\begin{verbatim}
{"ok": false}
\end{verbatim}

\begin{remark}[原因]
上述约束对应典型的非正则语言（需要栈/递归能力）。按第\ref{subsec:contract}小节合同，应直接判定失败而非给出“看似能用”的脆弱正则。
\end{remark}

\clearpage

%%%%%%%%%%%%%%%%%%%%%%%%%%%%%%%%%%%%%%%%%%%%%%%%%%%%%%%%%%%%
\section{MDQ 分割与外包证书化：把语义函子作为 HACA-PACER 的可调用微分基底}
\label{sec:mdq-outsourcing-haca-pacer}
%%%%%%%%%%%%%%%%%%%%%%%%%%%%%%%%%%%%%%%%%%%%%%%%%%%%%%%%%%%%

\begin{remark}[核心陈述（最快、不训练的关键等价）]
把“语义函子 = 生成可检验的标准化判别器（\code{ok:true}+标准形式 / \code{ok:false}）”这条工程路线，
可直接解释为“MDQ（微分动力量子）的分割与外包证书化”，并把这些证书作为 HACA-PACER 的可调用基础件。
其中：速度来自 LLM 的候选生成；可信来自确定性校验器的裁决与回归测试；复用来自证书化与版本化。
\end{remark}

本章将上述等价关系写成可执行的构造链：先把语义函子严格化为部分函子，
再把 MDQ 分割写成“局部结构变化单元”的序列输出，
最后把每个单元的判别器与标准化规则封装为可外包复验的证书，并对齐到 HACA 的五层结构与 PACER 的确定性驱动闭环。

\subsection{形式化等价：语义函子（部分函子）与 MDQ 分割}
\label{subsec:mdq-equivalence}

\begin{definition}[语义函子（一般化：标准化判别器）]
在第\ref{subsec:contract}小节中，我们将“通过时返回标准正则”作为值域。
为容纳更一般的可判别结构（例如安全有限自动机/小解析器/约束集合），本章引入一个抽象的标准化判别器空间 $\mathsf{CanonDet}$，
并把语义函子写成部分函子：
\[
F:\mathsf{Expr}\to \{\bot\}\sqcup \mathsf{CanonDet}.
\]
约定 $\bot$ 对应 \code{ok:false}；而 $\mathsf{CanonDet}$ 中的元素对应 \code{ok:true} 且携带一个可执行、可规范化的判别器负载。
\end{definition}

\begin{remark}[从正则到三类 witness（仍然可验证、且不需要训练）]
为了降低 $\bot$ 的比例，同时保持确定性可复验，工程上可将 $\mathsf{CanonDet}$ 的 witness 设计为三档（由快到强）：
\begin{enumerate}
  \item 安全正则（\textsf{RE2} 子集，见第\ref{subsec:contract}小节）；
  \item 有限自动机/约束组合（仍属正则语言，但实现上更可控）；
  \item 小解析器（限定语法、限定深度、可证明终止）。
\end{enumerate}
无论 witness 选哪一档，对外合同仍保持为“通过 + 标准形式 / 失败”。
\end{remark}

\begin{definition}[MDQ 分割算子与 MDQ 单元]
令 $\mathsf{MDQ}$ 表示“微分动力量子”（Micro-Differential Quantum）的结构空间。
MDQ 分割的目标是：把输入表达分解成一串“最小可验证的局部结构变化单元”。
形式上给出分割算子
\[
\mathsf{MDQSeg}:\mathsf{Expr}\to (\mathsf{MDQ})^{*},
\]
并约定每个输出单元可写为四元组
\[
\mathsf{MDQ}_i=(\mathrm{span}_i,\ \mathrm{type}_i,\ \mathrm{det}_i,\ \mathrm{canon}_i),
\]
其中 $\mathrm{span}_i$ 是局部作用范围，$\mathrm{type}_i$ 是类型标签，
$\mathrm{det}_i$ 是局部判别器（witness），$\mathrm{canon}_i$ 是与之配套的标准化重写（canonicalization）。
\end{definition}

\begin{remark}[等价点（工程实现）]
“最快方式、不训练”的等价点可被落实为三条硬约束：
\begin{enumerate}
  \item LLM 只负责为 $\mathrm{det}_i$ 与 $\mathrm{canon}_i$ 提供候选（提案）；
  \item 确定性校验器裁决候选是否可接受，并在通过时将其固化为标准形式；
  \item 一旦通过，即签发证书并进入可复用库，供 HACA-PACER 反复调用（外包）。
\end{enumerate}
\end{remark}

\begin{proposition}[“标准判别器”作为 MDQ 的微分基底（接口化表述）]
设 $\mathcal{B}\subseteq \mathsf{CanonDet}$ 为一组已证书化、可调用的标准判别器集合。
将表达的变化投影到该集合上，可视为对 $\mathcal{B}$ 的“坐标化”：
\[
\Pi_{\mathcal{B}}:\mathsf{Expr}\to \mathcal{B}^{*},
\]
其中 $\Pi_{\mathcal{B}}(x)$ 输出的是一串可复验的局部判别器调用与标准化重写（即 MDQ 单元的可执行载荷）。
在该口径下：
\begin{itemize}
  \item 语义函子 $F$ 输出的“标准判别器”可被理解为 MDQ 的微分基底元；
  \item MDQ 分割 $\mathsf{MDQSeg}$ 则是把输入表达的结构变化拆解并投影到这一组基底上的过程。
\end{itemize}
\end{proposition}

\subsection{证书外包：把判别器封装为可验证三元组}
\label{subsec:mdq-certificate}

\begin{definition}[可外包证书（Spec--Witness--Verifier）]
称一个 MDQ 判别器的外包证书为三元组
\[
\mathsf{Cert}=(\mathsf{Spec},\ \mathsf{Witness},\ \mathsf{Verify}),
\]
其中：
\begin{itemize}
  \item $\mathsf{Spec}$（声明）刻画该判别器所覆盖的语义边界；
  \item $\mathsf{Witness}$（见证）给出具体判别器（正则/自动机/小解析器）与标准化规则；
  \item $\mathsf{Verify}$（验证器）是一段确定性程序，用于检查 witness 是否满足 spec，并输出可审计结果。
\end{itemize}
\end{definition}

工程上可将证书回包设计为：外部只认 \code{ok} 与 \code{payload}，其余字段用于回归、审计与版本治理。

\paragraph{通过（示例）：}
\begin{verbatim}
{
  "ok": true,
  "mdq_type": "UUID",
  "payload": {
    "detector_kind": "regex",
    "detector": "^[0-9a-fA-F]{8}-...$",
    "canonicalize": "lowercase"
  },
  "test_digest": "sha256:...",
  "version": "mdq.uuid.v1"
}
\end{verbatim}

\paragraph{失败（示例）：}
\begin{verbatim}
{"ok": false, "mdq_type": "UUID"}
\end{verbatim}

\begin{remark}[外包的含义]
外包的含义是：HACA-PACER 不需要知道 witness 从何而来（LLM 提案、人工编写或迁移复用），
它只信任 $\mathsf{Verify}$ 能在本地或可信服务中对证书进行确定性复验并给出一致结论。
\end{remark}

\subsection{“最快、不训练”的闭环：LLM 提案，系统裁决，证书入库}
\label{subsec:fast-loop-no-training}

速度来自 LLM 的候选生成；可信来自裁决器；复用来自证书库。一个最小闭环为：
\begin{enumerate}
  \item \textbf{LLM 提案：}输出候选 witness（正则/自动机/小解析器）与候选标准化规则，并给出正例/反例与失败条件；
  \item \textbf{校验器裁决：}执行白名单、编译、例子验证、随机与对抗测试、复杂度上限等（参见第\ref{subsec:two-stage-pipeline}与第\ref{subsec:metamorphic-testing}
    小节）；
  \item \textbf{签发证书：}通过则版本化入库，并固定回归用的 \code{test_digest}；不通过则返回失败并记录原因聚类。
\end{enumerate}

\begin{remark}[为何不应把 witness 固定为正则]
正则是最快的一类 witness，但 MDQ 分割迟早会遇到“正则不够表达”的结构（嵌套、配对、跨段依赖等）。
把 witness 设计为“正则/自动机/小解析器”三档，可显著减少 \code{ok:false}，且仍保持可验证与不训练。
\end{remark}

\subsection{与 HACA 的五层结构对齐：MDQ 类型分层与状态表达}
\label{subsec:align-haca}

HACA 的（词法/语法/句法/文法/语义）五层结构可直接作为 MDQ 类型分层的口径：
\begin{itemize}
  \item $\mathsf{MDQ}^{(1)}$：词法量子（token/字符类/编码与空白/标点）；
  \item $\mathsf{MDQ}^{(2)}$：语法量子（局部结构：短语、限定词、数量、比较）；
  \item $\mathsf{MDQ}^{(3)}$：句法量子（依存关系、主谓宾框架的可判别片段）；
  \item $\mathsf{MDQ}^{(4)}$：文法量子（风格/格式约束：模板、字段协议、接口契约）；
  \item $\mathsf{MDQ}^{(5)}$：语义量子（实体类型、单位制、时空约束、域内本体）。
\end{itemize}

在证书化口径下，可将输入表达 $x$ 的“分层状态”写为：
\[
\mathsf{State}(x)=\bigl(\Pi_1(x),\Pi_2(x),\Pi_3(x),\Pi_4(x),\Pi_5(x)\bigr),
\]
其中 $\Pi_k$ 由 $\mathsf{MDQSeg}$ 的聚合与标准化产生，并以证书链作为可审计解释。

\subsection{PACER：规划--调用--检查--导出--修复（确定性驱动）}
\label{subsec:pacer}

在“不训练”约束下，PACER 更像一个确定性的“规划--检查--修复”驱动器：
\begin{itemize}
  \item P（Plan）：基于当前分层状态与目标接口契约，规划下一步需要补齐/修复的 MDQ；
  \item A（Act）：调用对应的 MDQ 证书（或触发新证书生成闭环）；
  \item C（Check）：检查是否消除了歧义/冲突（可将其记账为“表达上同调障碍”的受控版本）；
  \item E（Explain/Export）：输出可解释的标准化表达与证书链；
  \item R（Repair）：若失败，则提升到更高层 witness（或改写输入以回到可验证边界）。
\end{itemize}

\begin{remark}[用“新增边缘算子（新证书）”杀掉障碍]
当检查阶段出现不可控歧义/冲突时，PACER 的动作不是“继续猜”，而是引入新的可证书化判别器与标准化重写，
把障碍逼到系统设定的边界之外，从而将输入拉回可验证的标准表达域。
\end{remark}

\subsection{起步基底：一组高频 MDQ 证书（建议）}
\label{subsec:bootstrap-basis}

为了尽快让 HACA-PACER 获得可调用的“微分基底”，可先选一组小而高频的 MDQ 证书集合：
\begin{itemize}
  \item 数字/单位/区间（含货币、百分比、时间范围）；
  \item UUID/邮箱/URL/IP/日期时间；
  \item API 字段名与允许字符集（字段协议）；
  \item 常见自然语言约束词：至少/至多/介于/不超过/必须/可选。
\end{itemize}

\begin{remark}[扩展策略]
一旦起步基底建立起来，后续扩展主要是“补基底 + 回归测试 + 版本治理”，而非训练；
系统的可控性由证书复验与鲁棒性回归给出，系统的覆盖面由证书库规模给出。
\end{remark}

\clearpage

%%%%%%%%%%%%%%%%%%%%%%%%%%%%%%%%%%%%%%%%%%%%%%%%%%%%%%%%%%%%
\section{MDQ 粒度审计：把“可信”降解为可复验、可追溯、可回放的最小单元}
\label{sec:mdq-granular-audit}
%%%%%%%%%%%%%%%%%%%%%%%%%%%%%%%%%%%%%%%%%%%%%%%%%%%%%%%%%%%%

\begin{remark}[核心陈述（稳妥的起步路径）]
先把审计做到 MDQ（微分动力量子）这一层，等价于把“系统是否可信”的问题，
从“整段自然语言是否可信”降解为“每个可验证的微小结构单元是否可复验、可追溯、可回放”。
其原因在于：MDQ 的产物本质是“证书化的判别器 + 标准化重写”，天然具备审计所需的三要素——可记录、可复验、可对比。
\end{remark}

本章给出一套不依赖训练的 MDQ 粒度审计方案：以追加写、不可变的事件流承载审计证据，
并以三条主线组织审计能力（证书链审计、决策可回放审计、差分审计），从而为后续 HACA-PACER 的每一步调用提供天然的证书链与回归基线。

\subsection{审计对象：MDQ 片段、证书指纹与策略快照}
\label{subsec:audit-objects}

\begin{definition}[MDQ 审计事件（最小可复验记录）]
一次请求在 MDQ 粒度的审计记录可抽象为事件
\[
e=\bigl(\mathrm{req},\ \mathrm{ts},\ \mathrm{pipe},\ d_{\mathrm{in}},\ \mathsf{Seg},\ d_{\mathrm{out}},\ \mathsf{Policy}
  \bigr),
\]
其中：
\begin{itemize}
  \item $\mathrm{req}$ 为 \code{request\_id}，$\mathrm{ts}$ 为时间戳；
  \item $\mathrm{pipe}$ 为流水线版本（\code{pipeline_version}）；
  \item $d_{\mathrm{in}}$ 为输入摘要（\code{input_digest}，对原始输入做哈希；原文可选采样留存并脱敏/加密）；
  \item $\mathsf{Seg}$ 为 MDQ 分割结果列表（见下定义）；
  \item $d_{\mathrm{out}}$ 为输出摘要（\code{output_digest}，对标准化表达/标准正则等做哈希）；
  \item $\mathsf{Policy}$ 为策略快照（\code{policy_snapshot}：白名单规则版本、复杂度阈值版本等）。
\end{itemize}
\end{definition}

\begin{definition}[MDQ 分割条目（segment）]
将一次 MDQ 分割得到的条目写为
\[
s=\bigl(\mathrm{mdq\_type},\ d_{\mathrm{span}},\ \mathrm{cert\_id},\ \mathsf{Hash},\ \mathrm{ver},\ \mathrm{ok},\
  \mathrm{reason}\bigr),
\]
其中 $\mathrm{mdq\_type}$ 为类型，$d_{\mathrm{span}}$ 为片段摘要（\code{span_digest}），
$\mathrm{cert\_id}$ 与 $\mathrm{ver}$ 表示证书 ID/版本（如 \code{mdq.uuid.v1}），
$\mathsf{Hash}$ 为证书指纹（\code{spec_hash}、\code{witness_hash}、\code{verifier_version}、\code{test_digest}），
$\mathrm{ok}$ 为验证结果，$\mathrm{reason}$ 为失败原因码（白名单违规、编译失败、对抗测试失败、复杂度超阈值等）。
\end{definition}

\begin{remark}[事件流形态]
推荐将审计日志实现为“追加写、不可变”的事件流：一旦写入，任何更改都必须以新事件形式追加并留下链式指纹，
从而支持追责、对比与回放；同时将隐私风险控制在“digest-first”的最小留存策略之下（见第\ref{subsec:privacy-min-retention}小节）。
\end{remark}

\subsection{审计主线 I：证书链审计（可追溯）}
\label{subsec:certificate-chain-audit}

证书链审计回答审计的核心问题：\emph{这次判定依赖了哪些证书？这些证书当时是否通过确定性验证？之后是否被替换或升级？}

建议至少记录以下信息（均可由第\ref{subsec:audit-objects}小节的事件结构覆盖）：
\begin{itemize}
  \item 分割得到的 MDQ 列表：\code{mdq_type}、\code{span_digest}、顺序与层级；
  \item 每个 MDQ 绑定的证书 ID/版本：\code{cert_id}；
  \item 证书指纹：\code{spec_hash}、\code{witness_hash}、\code{verifier_version}、\code{test_digest}；
  \item 验证结果与失败原因码：\code{ok}、\code{reason_code}。
\end{itemize}

\begin{remark}[可追溯性的判据]
若任意一次输出都能在事件流中回溯到一条有限的证书链，并能以相同的 \code{verifier_version} 复验出相同的 \code{ok/reason_code}，
则该输出在证书层面是可追溯的；反之若证书版本、验证器版本或策略快照缺失，则追溯断裂，审计结论不成立。
\end{remark}

\subsection{审计主线 II：决策可回放审计（可复验）}
\label{subsec:replay-audit}

\begin{definition}[提案--裁决分离记录]
将一次判定拆分为两类可审计产物：
\begin{enumerate}
  \item LLM 提案：候选判别器（如正则）、正反例、模型与提示词版本；
  \item 确定性裁决：校验器输出（通过/拒绝）与拒绝原因码，以及通过时的规范化结果与证书指纹。
\end{enumerate}
审计口径规定：\emph{只信任确定性裁决的输出}；LLM 提案仅作为可追责的输入材料与回归样例来源。
\end{definition}

\begin{definition}[回放算子与一致性]
给定一条历史事件 $e$，记其裁决所用验证器版本为 $\mathrm{ver}(e)$，策略快照为 $\mathsf{Policy}(e)$。
用相同版本验证器与相同策略快照对同一输入摘要（或其受控可复原输入）执行复验，称为回放（replay）。
若回放得到的 MDQ 级结果（每个条目的 \code{ok/reason_code} 与规范化输出摘要）与原事件一致，则称该事件满足回放一致性。
\end{definition}

\begin{remark}[漂移与篡改检测]
回放一致性使得“事后重放”成为可能：抽样选取历史输入，拉起当时的 \code{verifier_version} 与 \code{policy_snapshot} 复验，
即可检测语义漂移、验证器变更、策略篡改或证书库污染等问题。
\end{remark}

\subsection{审计主线 III：差分审计（可对比、可定位）}
\label{subsec:differential-audit}

当升级某类 MDQ 证书（例如从 \code{mdq.uuid.v1} 到 \code{mdq.uuid.v2}）时，建议做 MDQ 级差分，而非只看“整条输入通过率”。
差分审计关注可定位的指标：
\begin{itemize}
  \item 哪些 \code{mdq_type} 的通过率变化最大；
  \item 哪些标准判别器发生了结构变化（规范化后仍不同）；
  \item 行为等价是否保持（同一批测试串上 \code{match} 结果是否一致）；
  \item 漂移集中在哪些语义表达变体上（同义改写、噪声扰动、对抗样例）。
\end{itemize}

\begin{definition}[MDQ 差分报告（摘要口径）]
对固定回归集与固定策略快照，比较版本 $\mathrm{v1}$ 与 $\mathrm{v2}$ 的事件集合，输出按 \code{mdq_type} 聚合的差分摘要：
通过率变化、典型失败原因码分布、规范化输出变化摘要，以及行为等价抽样检验统计。
\end{definition}

\begin{remark}[定位粒度]
差分审计能把问题定位到“某一个 \code{mdq\_type} + 某一个证书版本 + 某一类输入形态”，显著降低排查成本。
\end{remark}

\subsection{事件日志建议字段（最小但可复验）}
\label{subsec:log-schema-minimal}

工程上建议每次调用至少落以下字段（尽量少，但足够复验）：
\begin{itemize}
  \item \code{request_id}、\code{timestamp}、\code{pipeline_version}；
  \item \code{input_digest}（必要时额外保存脱敏/加密原文的引用指针）；
  \item \code{mdq_segments}：\code{[(mdq_type, span_digest, cert_id, cert_hashes, verifier_version, ok, reason_code)]}；
  \item \code{output_digest}；
  \item \code{policy_snapshot}（白名单规则版本、复杂度阈值版本等）。
\end{itemize}

\subsection{隐私与安全：digest-first 的最小留存策略}
\label{subsec:privacy-min-retention}

\begin{remark}[默认只存摘要]
隐私与安全上，可默认只存 digest，不存原文；需要可解释性时再对少量样本做脱敏留存或加密留存。
这样既能审计（复验、对比、回放），又能降低敏感数据长期留存风险。
\end{remark}

\subsection{上线门槛：两条硬指标（建议）}
\label{subsec:audit-gates}

若要将 MDQ 粒度审计作为“可上线”的模块，建议至少设置两条门槛指标：
\begin{itemize}
  \item \textbf{可回放一致性：}抽样重放历史请求，MDQ 级结果一致率接近 $100\%$（允许极少数来自外部依赖的例外，但必须可解释并可审计）；
  \item \textbf{证书漂移可解释性：}任何 \code{mdq_type} 的版本升级，必须产出一份 MDQ 级差分报告（通过率变化、典型失败样例、规范化变化摘要、行为等价统计）。
\end{itemize}

\clearpage

%%%%%%%%%%%%%%%%%%%%%%%%%%%%%%%%%%%%%%%%%%%%%%%%%%%%%%%%%%%%
\section{MDQ 作为“可移植层”：同伦类归属的提案--复验拆分与留痕审计}
\label{sec:mdq-portable-layer-homotopy-membership}
%%%%%%%%%%%%%%%%%%%%%%%%%%%%%%%%%%%%%%%%%%%%%%%%%%%%%%%%%%%%

\begin{remark}[核心陈述（可移植层）]
把 MDQ 做成“可移植层”的关键，是把“同伦类归属的判定”拆成两部分：
\begin{itemize}
  \item LLM 只负责提出提案：它认为输入属于哪个表达同伦类/对应哪个标准谓词；
  \item MDQ 负责用确定性的标准谓词复验，并把偏差留痕（可追溯、可回放、可对比）。
\end{itemize}
于是更换任何 LLM（不同厂商、不同版本、不同上下文长度）都不会破坏审计闭环；
差异只体现在提案质量与速度上，而不会影响最终裁决的可复验性。
\end{remark}

本章用最少符号把上述逻辑写成可执行的接口，并给出 MDQ 粒度需要落盘的审计字段（与第\ref{sec:mdq-granular-audit}节一致且更聚焦于“同伦类归属”问题）。

\subsection{表达同伦类与标准谓词：从 \texorpdfstring{$[E]$}{[E]} 到可执行判别器}
\label{subsec:homotopy-class-predicate}

\begin{definition}[表达同伦类与标准谓词]
令 $\mathsf{Expr}$ 为表达域，并令 $\mathsf{H}$ 表示其“表达同伦类”的集合（或选定口径下的等价类集合）。
对每个类 $h\in\mathsf{H}$（也记作 $[E]$），固定绑定一个标准逻辑谓词 $P_h$，并要求其可落为一个可执行判别器：
\[
r_h\in \mathsf{CanonDet},
\]
其中 $\mathsf{CanonDet}$ 见第\ref{subsec:mdq-equivalence}小节。
在“正则起步”口径下，可取 $r_h$ 为安全正则（\textsf{RE2} 子集）或与之等价的有限自动机。
\end{definition}

\begin{definition}[标准化函数]
给定一个标准化函数
\[
C:\mathsf{Expr}\to \mathsf{Expr}_c,
\]
其中 $\mathsf{Expr}_c\subseteq \mathsf{Expr}$ 表示标准化后的表达空间（规范形空间），$C$ 将输入归一到标准形（例如统一空白、大小写、全半角、符号别名等）。
\end{definition}

\begin{definition}[确定性归属判定（指示函数）]
记 $L(r)$ 为判别器 $r$ 接受的语言（其可判定字符串集合）。
对任意 $h\in\mathsf{H}$ 与 $x\in\mathsf{Expr}$，定义确定性归属判定：
\[
\mathrm{In}_h(x):=\mathbf{1}\bigl[C(x)\in L(r_h)\bigr]\in\{0,1\}.
\]
\end{definition}

\begin{remark}[工程解释]
式中 $C(x)$ 把输入先拉回标准形，再由 $r_h$ 判定其是否落入类 $h$ 的可接受域。
因此“同伦类 $\leftrightarrow$ 标准谓词/判别器”的体系一旦证书化（见第\ref{subsec:mdq-certificate}小节），
归属判定就不再依赖 LLM 的口头解释，而只依赖可复验的确定性谓词。
\end{remark}

\subsection{LLM 提案与 MDQ 复验：裁决以确定性为准}
\label{subsec:proposal-vs-verification}

\begin{definition}[LLM 提案（proposal）]
LLM 对输入 $x$ 给出一个提案对
\[
(\hat{h}(x),\hat{y}(x))\in \mathsf{H}\times\{0,1\},
\]
其中 $\hat{h}(x)$ 表示其认为 $x$ 归属的同伦类，$\hat{y}(x)$ 表示其对“属于/不属于”的布尔判断。
（实现中也可允许返回候选类集合 $\widehat{H}(x)\subseteq\mathsf{H}$，再由 MDQ 逐个复验。）
\end{definition}

\begin{definition}[MDQ 复验（verification）]
给定 LLM 的提案 $\hat{h}(x)$，MDQ 侧计算确定性结果
\[
y^\ast(x):=\mathrm{In}_{\hat{h}(x)}(x).
\]
系统的最终裁决以 $y^\ast(x)$ 为准，并将 $(\hat{h}(x),\hat{y}(x))$ 仅作为可审计的提案元数据保留。
\end{definition}

\begin{remark}[可移植性的来源]
由于最终裁决由 $\mathrm{In}_h$（确定性谓词）给出，LLM 只提供提案；
因此更换模型只会改变 $\hat{h},\hat{y}$ 的分布与质量，不会改变裁决机制本身的可复验性。
这正是“MDQ 作为可移植层”的含义。
\end{remark}

\subsection{分歧类型：误识别与漏识别的最小逻辑刻画}
\label{subsec:discrepancy-types}

\begin{definition}[分歧事件与分歧类型]
对输入 $x$，若 $\hat{y}(x)\neq y^\ast(x)$，称发生分歧事件。
分歧可分为两类：
\begin{enumerate}
  \item \textbf{误识别（不属于却判成属于）.}
  \[
  \hat{y}(x)=1\ \wedge\ y^\ast(x)=0;
  \]
  \item \textbf{漏识别（属于却判成不属于）.}
  \[
  \hat{y}(x)=0\ \wedge\ y^\ast(x)=1.
  \]
\end{enumerate}
\end{definition}

\begin{remark}[留痕即审计]
只要把上述两类分歧写入审计日志（并带上证书指纹与策略快照），就完成了“留痕即审计”：
换一个模型跑同一批输入，分歧比例、分布以及集中在哪些同伦类上，都可直接量化为“提案质量与容错率由 LLM 决定”。
\end{remark}

\subsection{MDQ 粒度审计字段：提案、标准化、谓词指纹与复验结果}
\label{subsec:homotopy-audit-fields}

为使审计可复验且可对比，建议在 MDQ 粒度至少记录以下字段（轻量但足够回放）：
\begin{itemize}
  \item \code{input_digest}：原始输入哈希（必要时再存脱敏原文或其加密引用）；
  \item \code{canon_digest}：标准化后 $C(x)$ 的哈希，以及标准化步骤摘要（便于解释“为何判定改变”）；
  \item \code{llm_proposal}：$\hat{h},\hat{y}$、模型版本与 \code{prompt_version}；
  \item \code{predicate_id}：被调用的标准谓词/判别器 $r_h$ 的 \code{predicate_version} 与哈希（防止事后被悄然替换）；
  \item \code{verify_result}：确定性验证结果 $y^\ast$；
  \item \code{discrepancy}：是否分歧、分歧类型（第\ref{subsec:discrepancy-types}小节两类之一）与 \code{reason_code}
  （如白名单不通过、对抗测试失败、复杂度超阈值、或“仅差一步标准化”等）。
\end{itemize}

\begin{remark}[与证书外包的对齐]
若要把“同伦类 $\leftrightarrow$ 标准谓词/判别器”的体系做成证书外包，
建议为每个同伦类引入稳定的 \code{class_id}，并规定：系统最终输出以 MDQ 复验为准；
同时保留 LLM 原始提案作为日志字段。这样审计既不丢信息，也不会让系统被 LLM 的偶发漂移拖着走。
\end{remark}

\clearpage

%%%%%%%%%%%%%%%%%%%%%%%%%%%%%%%%%%%%%%%%%%%%%%%%%%%%%%%%%%%%
\section{多次扫描的收敛：候选同伦类集合的单调收缩与固定点判据}
\label{sec:multi-scan-convergence}
%%%%%%%%%%%%%%%%%%%%%%%%%%%%%%%%%%%%%%%%%%%%%%%%%%%%%%%%%%%%

\begin{remark}[核心陈述（收敛过程）]
“多次扫描后语义就明确”可以理解为一个可复验的收敛过程：每一次扫描都把不确定性往外挤，
把表达逐步压到某个同伦类所对应的标准逻辑谓词（标准正则/标准判别器）上，直到结果稳定为止。
更具体地，若把“语义”视为一个候选同伦类集合，那么每扫一遍就可将候选集合收缩一次，
并最终达到一个不再变化的固定点。
\end{remark}

本章在第\ref{sec:mdq-portable-layer-homotopy-membership}节“提案--复验拆分”的基础上，
把多轮扫描写成一个单调迭代过程，并给出三类终止判据（单点、标准化稳定、分层稳定），以及证书基底不足时的两种残余收敛状态（歧义固定点与空集合）。

\subsection{候选集合与扫描算子：单调收缩的形式化}
\label{subsec:candidate-set-scan-operator}

\begin{definition}[候选同伦类集合]
令 $\mathsf{H}_{\mathrm{lib}}\subseteq \mathsf{H}$ 为当前证书库可覆盖的同伦类集合（即存在可复验谓词 $P_h$ 与判别器 $r_h$ 的类）。
对固定输入 $x\in\mathsf{Expr}$，定义第 $t$ 轮扫描后的候选集合为
\[
\mathsf{Cand}_t(x)\subseteq \mathsf{H}_{\mathrm{lib}}.
\]
一种最保守的初始化是 $\mathsf{Cand}_0(x):=\mathsf{H}_{\mathrm{lib}}$（先允许一切，再逐步排除）；也可用轻量规则先粗筛得到较小初值。
\end{definition}

\begin{definition}[第 $t$ 轮提案与可证伪剔除]
第 $t$ 轮中，LLM 给出候选类集合 $\widehat{H}_t(x)\subseteq \mathsf{Cand}_t(x)$（可由多模型/多提示聚合得到），
MDQ 对其中每个 $h\in \widehat{H}_t(x)$ 计算确定性复验值 $\mathrm{In}_h(x)$（见第\ref{subsec:homotopy-class-predicate}小节）。
定义第 $t$ 轮被证伪的候选集合为
\[
\mathsf{Elim}_t(x):=\{h\in \widehat{H}_t(x):\mathrm{In}_h(x)=0\}.
\]
并定义一次扫描更新为
\[
\mathsf{Cand}_{t+1}(x):=\mathsf{Cand}_t(x)\setminus \mathsf{Elim}_t(x).
\]
\end{definition}

\begin{proposition}[单调性（候选集合不增）]
对任意 $t\ge 0$，有
\[
\mathsf{Cand}_{t+1}(x)\subseteq \mathsf{Cand}_t(x).
\]
\end{proposition}

\begin{remark}[速度与完备性的权衡]
若每轮只验证少量 $h$（即 $\widehat{H}_t(x)$ 很小），则每轮剔除有限，但成本低；
若接近穷举验证 $\mathsf{Cand}_t(x)$，则收缩更快但成本高。
工程上通常选“提案引导的稀疏验证”，并用差分审计与回归驱动提案质量提升（见第\ref{subsec:differential-audit}小节）。
\end{remark}

\subsection{固定点与终止条件：语义明确的三类判据}
\label{subsec:fixed-point-termination}

\begin{definition}[固定点]
若存在 $T$ 使得
\[
\mathsf{Cand}_{T+1}(x)=\mathsf{Cand}_T(x),
\]
则称 $\mathsf{Cand}_T(x)$ 为在当前证书基底与策略快照下的固定点候选集合。
\end{definition}

\begin{remark}[三类终止判据（任一满足即可停机）]
在工程口径下，通常采用如下三类终止判据（满足其一即可宣告“语义在当前基底下已明确/已封闭”）：
\begin{enumerate}
  \item \textbf{单点判据：}候选集合收缩为单点，即 $|\mathsf{Cand}_T(x)|=1$。
  此时唯一剩余的 $h^\ast\in\mathsf{Cand}_T(x)$ 的谓词 $P_{h^\ast}$（以及判别器 $r_{h^\ast}$）就是标准语义锚点。
  \item \textbf{标准化稳定判据：}连续两轮或三轮扫描后，标准化输出不再变化，且已启用的 MDQ 证书均复验通过。
  该判据表达“再怎么扫也不会产生新的结构信息”，语义在当前证书基底下已经封闭。
  \item \textbf{分层稳定判据：}从 $\mathsf{MDQ}^{(1)}$ 到 $\mathsf{MDQ}^{(5)}$ 的分层输出均稳定（见下小节），
  等价于“语义从形式到含义的各层都锁定了”。
\end{enumerate}
\end{remark}

\subsection{分层稳定：\texorpdfstring{$\mathsf{MDQ}^{(1)}\sim \mathsf{MDQ}^{(5)}$}{MDQ(1)\textasciitilde MDQ(5)} 的稳定通过}
\label{subsec:layerwise-stability}

沿用第\ref{subsec:align-haca}小节的五层分解，记 $\Pi_k(x)$ 为第 $k$ 层的聚合与标准化输出（由 $\mathsf{MDQSeg}$ 产生）。
将“分层稳定”具体化为：连续多轮扫描后，下列现象均不再发生（或其摘要不再变化），且对应证书复验持续通过：
\begin{itemize}
  \item $\mathsf{MDQ}^{(1)}$：词法层不再出现新的 token 切分或编码/空白差异；
  \item $\mathsf{MDQ}^{(2)}$：局部结构（数量、比较、限定）不再摇摆；
  \item $\mathsf{MDQ}^{(3)}$：句法可判别片段稳定；
  \item $\mathsf{MDQ}^{(4)}$：模板/协议/契约匹配稳定；
  \item $\mathsf{MDQ}^{(5)}$：实体类型、单位制、时空约束、域内本体映射稳定。
\end{itemize}

\begin{remark}[可审计实现]
分层稳定可用“摘要一致性”实现：对每层输出保存 \code{digest}（以及少量可解释摘要），
要求连续两轮或三轮的 \code{digest} 一致，并且该层涉及的证书指纹（\code{predicate_id}/\code{cert_id} 等）未变且复验通过。
这与第\ref{sec:mdq-granular-audit}节的事件流审计完全兼容。
\end{remark}

\subsection{留痕与收敛：误收/漏收在多轮扫描中的角色}
\label{subsec:discrepancy-and-convergence}

多次扫描并不是让 LLM 自己“变得更可靠”，而是用确定性 MDQ 判别器把两类错误显式化并持续留痕（第\ref{subsec:discrepancy-types}小节）：
\begin{itemize}
  \item \textbf{误收（不属于却判为属于）}会被复验当场拦下并记为分歧事件，形成“哪些输入形态容易被某个 LLM 误判”的清单；
  \item \textbf{漏收（属于却判为不属于）}可通过“多模型/多提示/多轮提案”提高召回率：
  只要某一轮把正确的同伦类 $h$ 提出来，MDQ 就能复验成立，从而把语义压到正确的标准谓词上。
\end{itemize}

\begin{remark}[收敛的工程含义]
在该口径下，扫描轮数增加带来的收益主要体现在“漏收”的缓解（更高的提案覆盖），而“误收”始终由确定性复验兜底。
因此系统稳定性由 MDQ 裁决层给出，LLM 主要影响效率与到达固定点所需的轮数。
\end{remark}

\subsection{证书基底覆盖与残余收敛：歧义固定点与空集合}
\label{subsec:basis-coverage-residual}

\begin{remark}[前提：基底覆盖足够]
“多次扫描能明确语义”的前提，是当前 MDQ 证书基底对目标领域的覆盖足够。
否则迭代同样会趋于稳定，但可能稳定在两种残余状态之一；重要的是，这两种状态都会被审计日志明确暴露为可定位的工程问题。
\end{remark}

\begin{definition}[两类残余状态]
设迭代达到固定点 $\mathsf{Cand}_\infty(x)$（可取某个足够大的 $T$ 的 $\mathsf{Cand}_T(x)$ 作为近似）。
\begin{enumerate}
  \item \textbf{歧义固定点：}$|\mathsf{Cand}_\infty(x)|>1$。
  表示在现有证书库下仍然歧义（多个同伦类均未被证伪或无法在成本约束下分离）。
  \item \textbf{空集合：}$\mathsf{Cand}_\infty(x)=\varnothing$。
  表示输入与当前规则体系冲突，或缺少必要的谓词/判别器/标准化规则，使得所有可覆盖同伦类均不可接受。
\end{enumerate}
\end{definition}

\begin{remark}[工程闭环]
两类残余状态并非坏事：它们把“缺基底/缺规则”转化为可追溯的差分事件，
从而提示应补齐哪一类 MDQ、哪一条谓词、或哪一段标准化规则，并可直接纳入差分审计与回归集构造（第\ref{subsec:differential-audit}小节）。
\end{remark}

\subsection{接入 HACA-PACER：固定点迭代的控制环}
\label{subsec:haca-pacer-fixed-point-loop}

若将多次扫描放入 HACA-PACER，则它自然成为一个固定点迭代的控制环：每一轮优先扫描并补齐那些最不稳定、最常产生分歧的 MDQ 类型，
直到满足第\ref{subsec:fixed-point-termination}小节的终止判据为止。

\begin{definition}[不稳定性驱动的扫描选择（示意）]
可按审计事件流为每个 MDQ 类型（或每一层 $k$）定义不稳定性评分 $\mathrm{Instab}$（例如分歧率、拒绝率、失败原因码的加权和等）。
令第 $t$ 轮需要优先处理的类型为
\[
\tau_t:=\arg\max_{\tau}\ \mathrm{Instab}_t(\tau),
\]
并围绕 $\tau_t$ 执行“提案--复验--修复/补证书”的闭环。
\end{definition}

\begin{remark}[可复验的收敛判据]
上述控制环的价值在于：语义的“明确”不依赖直觉，而依赖可复验的收敛判据与可回放的审计证据。
这使得系统既可上线（有门槛与指标），也可持续迭代（靠差分报告与回归集推进基底覆盖）。
\end{remark}

\clearpage

%%%%%%%%%%%%%%%%%%%%%%%%%%%%%%%%%%%%%%%%%%%%%%%%%%%%%%%%%%%%
\section{同伦类指纹：语义哈希与证书哈希的双层锚定（防洗稿口径）}
\label{sec:homotopy-fingerprint-hash}
%%%%%%%%%%%%%%%%%%%%%%%%%%%%%%%%%%%%%%%%%%%%%%%%%%%%%%%%%%%%

\begin{remark}[核心陈述（从“LLM 一句话”到“稳定指纹”）]
一旦“收敛的对象”不再是某个 LLM 的一句话，而是一个同伦类的稳定哈希（同伦类指纹），
就可以把“语义归属”变成可复验、可追溯且难以事后篡改的对象。
在工程层面，这至少支持“防洗稿”一类需求：表层措辞改变，但若仍落在同一同伦类，其指纹保持不变；反之若更换谓词/证书实现，则会留下可检测差异。
\end{remark}

本章在第\ref{sec:multi-scan-convergence}节的收敛口径与第\ref{sec:mdq-granular-audit}节的事件流审计口径之上，
引入两层哈希：\emph{语义同伦类哈希}（跨版本稳定，尽量与实现无关）与\emph{证书实现哈希}（与版本绑定，用于回放与防篡改）。
并进一步给出分层/局部指纹用于细粒度对比。

\subsection{固定对应关系：同伦类 \texorpdfstring{$\leftrightarrow$}{<->} 标准谓词 \texorpdfstring{$\leftrightarrow$}{<->} 标准判别器}
\label{subsec:fingerprint-fixed-correspondence}

沿用第\ref{subsec:homotopy-class-predicate}小节的设置：每个表达同伦类 $h\in\mathsf{H}$ 绑定一个标准谓词 $P_h$，
并落成一个可复验的标准判别器 $r_h\in\mathsf{CanonDet}$；并存在标准化函数 $C(\cdot)$。

\begin{definition}[稳定类编号 \code{class_id}]
为每个同伦类 $h\in\mathsf{H}$ 指定一个稳定编号
\[
\mathrm{ClassID}:\mathsf{H}\to \mathsf{ID},
\]
其中 $\mathsf{ID}$ 可取字符串/UUID 等稳定编码（如“UUID 格式类”“日期时间类”“某领域术语模板类”等）。
要求 $\mathrm{ClassID}$ 的含义在版本演进中保持稳定：谓词实现可升级，但类编号不应随实现细节变化。
\end{definition}

\begin{definition}[收敛类（单点情形）]
若对输入 $x$ 的多轮扫描在第\ref{subsec:fixed-point-termination}小节的单点判据下终止，
即存在 $T$ 使得 $\mathsf{Cand}_T(x)=\{h^\ast\}$，
则定义其收敛类为
\[
[E(x)]:=h^\ast.
\]
\end{definition}

\begin{remark}[非单点时的处理]
若固定点候选集合仍大于 $1$（歧义固定点）或为空（冲突/缺基底），
则不应宣称“语义指纹已明确”；应返回失败或返回“候选集合指纹/冲突指纹”，并由审计闭环驱动补齐基底（第\ref{subsec:basis-coverage-residual}小节）。
\end{remark}

\subsection{语义同伦类哈希：跨版本稳定的 \texorpdfstring{$H_{\mathrm{sem}}$}{H_sem}}
\label{subsec:semantic-homotopy-hash}

\begin{definition}[语义同伦类哈希]
在 $[E(x)]$ 已定义的前提下，定义
\[
H_{\mathrm{sem}}(x):=\operatorname{Hash}\bigl(\mathrm{ClassID}([E(x)])\bigr),
\]
其中 $\operatorname{Hash}$ 为固定的摘要函数（例如 \textsf{SHA-256}\cite{fips180-4}），输入编码采用稳定序列化（例如 UTF-8 + 规范化）。
\end{definition}

\begin{remark}[语义稳定性（防洗稿的核心）]
只要两段表达在您定义的同伦等价（允许的改写、同义替换、格式扰动）下属于同一类，
则它们应收敛到同一 $[E]$，从而得到相同的 $H_{\mathrm{sem}}$。
这使得“仅靠润色/改格式”难以洗掉指纹；若要改变指纹，必须改变其同伦类归属（即改变语义锚点）。
\end{remark}

\subsection{证书实现哈希：可回放、可审计的 \texorpdfstring{$H_{\mathrm{cert}}$}{H_cert}}
\label{subsec:certificate-implementation-hash}

\begin{definition}[证书实现哈希（版本绑定）]
定义证书实现哈希为
\[
H_{\mathrm{cert}}(x):=\operatorname{Hash}\Bigl(
\mathrm{ClassID}([E(x)]),\
\mathrm{PredicateID}([E(x)]),\
\mathrm{WitnessHash}([E(x)]),\
\mathrm{VerifierVersion},\
\mathrm{CanonVersion}
\Bigr),
\]
其中：
\begin{itemize}
  \item $\mathrm{PredicateID}$ 为标准谓词/判别器的版本标识（如 \code{predicate_version}）；
  \item $\mathrm{WitnessHash}$ 为判别器本体的摘要（正则字符串/自动机结构/小解析器实现的哈希）；
  \item $\mathrm{VerifierVersion}$ 与 $\mathrm{CanonVersion}$ 分别为验证器与标准化器版本。
\end{itemize}
也可将 \code{policy_snapshot}（白名单/复杂度阈值版本）纳入哈希输入，以确保策略变更可被检测。
\end{definition}

\begin{remark}[篡改可检出（防伪/防赖账）]
任何事后“偷偷换判别器/换验证器/换标准化规则”的行为，都会改变上述版本字段或 witness 摘要，
从而导致 $H_{\mathrm{cert}}$ 变化。若再将 $H_{\mathrm{cert}}$ 写入追加写、不可变的链式日志（第\ref{subsec:audit-objects}小节），
则篡改将很难在不留痕的情况下完成。
\end{remark}

\subsection{收敛 \texorpdfstring{$\Rightarrow$}{=>} 指纹收敛：从 \texorpdfstring{$[E]$}{[E]} 到 \texorpdfstring{$H_{\mathrm{sem}}
  $}{H_sem}}
\label{subsec:convergence-implies-fingerprint}

\begin{proposition}[同伦类收敛蕴含语义指纹稳定]
若对输入 $x$ 的多轮扫描在单点判据下终止于 $[E(x)]$，
并且后续重复扫描不改变该收敛类，则 $H_{\mathrm{sem}}(x)$ 在该基底与策略快照下稳定不变。
\end{proposition}

\paragraph{说明.}
这是第\ref{sec:multi-scan-convergence}节“固定点收敛”的直接推论：$H_{\mathrm{sem}}$ 仅依赖稳定的 $\mathrm{ClassID}$ 与收敛类。

\subsection{分层/局部指纹：从全局锚点到“语义片段集合”}
\label{subsec:layer-local-fingerprint}

\begin{remark}[为何需要分层/局部指纹]
现实中的洗稿常见手法是“整体改写，但保留关键结构片段”。仅靠单个全局 $H_{\mathrm{sem}}$ 可能过粗；
而 MDQ 天然提供分层输出与局部片段，因此可以无训练地构造“指纹向量/指纹集合”，用于检测局部高重合。
\end{remark}

\begin{definition}[分层语义指纹向量（示意）]
对 HACA 五层（第\ref{subsec:align-haca}小节）分别定义层内语义指纹
\[
H^{(k)}_{\mathrm{sem}}(x):=\operatorname{Hash}\bigl(\mathrm{LayerID}_k(x)\bigr),
\qquad k=1,\dots,5,
\]
其中 $\mathrm{LayerID}_k(x)$ 表示第 $k$ 层在选定口径下的稳定标识（可由该层 MDQ 证书调用链与规范化输出摘要导出）。
进而定义指纹向量
\[
\vec{H}_{\mathrm{sem}}(x):=\bigl(H^{(1)}_{\mathrm{sem}}(x),\dots,H^{(5)}_{\mathrm{sem}}(x)\bigr).
\]
\end{definition}

\begin{definition}[语义片段指纹集合（示意）]
令第 $k$ 层产生的 MDQ 片段集合为 $\mathsf{Seg}^{(k)}(x)$。对每个片段 $s\in \mathsf{Seg}^{(k)}(x)$，
可定义局部指纹
\[
\operatorname{fp}(s):=\operatorname{Hash}\bigl(\mathrm{class\_id}(s),\ \mathrm{span\_digest}(s)\bigr),
\]
并令
\[
\mathsf{FP}^{(k)}(x):=\{\operatorname{fp}(s):s\in \mathsf{Seg}^{(k)}(x)\}.
\]
通过比较 $\mathsf{FP}^{(k)}(x)$ 的重合度（例如 Jaccard 相似度），可检测“局部结构语义”是否高度重叠。
\end{definition}

\subsection{边界与调参：同伦等价口径过宽/过窄的代价}
\label{subsec:fingerprint-boundary}

\begin{remark}[同伦等价生成规则是关键边界]
“防洗稿”能成立的前提，是您对“同伦等价”的生成规则（允许哪些改写、视为同一类）定义得贴近目标：
\begin{itemize}
  \item 规则过宽：不同内容可能落到同一同伦类，误报增加；
  \item 规则过窄：同一内容换一种表达可能落到不同类，漏报增加。
\end{itemize}
这恰好是 MDQ 的优势：可用审计日志里积累的“误收/漏收”样本精修边界，而不必训练模型（第\ref{subsec:discrepancy-and-convergence}小节）。
\end{remark}

\subsection{作为证书外包产物固化：输出 \texorpdfstring{$(H_{\mathrm{sem}},H_{\mathrm{cert}})$}{(H_sem,H_cert)} 并写入链式日志}
\label{subsec:fingerprint-as-certificate-output}

\begin{remark}[建议的最小输出]
若认可“同伦类指纹”作为正式产物，建议每次输出至少包含
\[
\bigl(H_{\mathrm{sem}}(x),\ H_{\mathrm{cert}}(x)\bigr),
\]
并把它们与 \code{class_id}/\code{predicate_id}/\code{verifier_version}/\code{canon_version} 一并写入不可变事件流。
这样它既是防洗稿的语义锚点，也是后续 HACA-PACER 推进时可调用的“语义证书接口”（可回放、可对比、可追责）。
\end{remark}

\clearpage

%%%%%%%%%%%%%%%%%%%%%%%%%%%%%%%%%%%%%%%%%%%%%%%%%%%%%%%%%%%%
\section{语义最小压缩与一致展开：\texorpdfstring{$\pi$}{pi}--\texorpdfstring{$G$}{G} 回环不变性与 HACA-PACER 最小闭环}
\label{sec:semantic-compress-expand-roundtrip}
%%%%%%%%%%%%%%%%%%%%%%%%%%%%%%%%%%%%%%%%%%%%%%%%%%%%%%%%%%%%

\begin{remark}[核心陈述（两句话）]
本章把“防洗稿/可控生成/可回放审计”等能力统一抽象为两句话：
\begin{enumerate}
  \item \textbf{语义最小压缩：}将任意风格/措辞/格式的表达 $x$ 投影到一个同伦类语义指纹（并附带少量必要参数），得到最小可复验的语义签名；
  \item \textbf{语义一致展开：}从语义签名出发，在不同风格约束下生成表述，并用同一套 MDQ 证书体系复验仍落在同一同伦类中，从而形成闭环。
\end{enumerate}
在工程意义上，这给出一个可运行的 HACA-PACER 最小闭环骨架：分层语义结构（HACA） + 规划/行动/检查/解释/修复（PACER），并且核心判定与收敛不依赖训练。
\end{remark}

\subsection{同伦等价与语义签名空间 \texorpdfstring{$\mathsf{SemSig}$}{SemSig}}
\label{subsec:semsig-space}

\begin{definition}[同伦等价关系（口径）]
在表达域 $\mathsf{Expr}$ 上固定一个等价关系 $\sim$（同伦等价）：
\[
x\sim y
\quad:\Longleftrightarrow\quad
\text{$x$ 与 $y$ 在允许的改写/风格变化/格式扰动下属于同一语义类。}
\]
该口径由 MDQ 的标准化规则、证书基底与判别器体系共同确定，并应通过审计与回归持续维护。
\end{definition}

\begin{definition}[语义签名空间]
定义语义签名空间 $\mathsf{SemSig}$ 为一类结构化数据的集合。典型地，一个签名为四元组
\[
s=\pi(x)=\bigl(\mathrm{ClassID},\ \mathrm{Param},\ H_{\mathrm{sem}},\ H_{\mathrm{cert}}\bigr),
\]
其中：
\begin{itemize}
  \item $\mathrm{ClassID}$ 是稳定类编号（见第\ref{subsec:fingerprint-fixed-correspondence}小节）；
  \item $\mathrm{Param}$ 是该类所需的最小参数集合（例如单位制、时区、字段协议参数等；没有则可为空）；
  \item $H_{\mathrm{sem}}$ 与 $H_{\mathrm{cert}}$ 分别为语义同伦类哈希与证书实现哈希（第\ref{subsec:semantic-homotopy-hash}--\ref{subsec:certificate-implementation-hash}小节）。
\end{itemize}
\end{definition}

\begin{remark}[“最小”的含义（工程口径）]
所谓“语义最小压缩”，并不要求在数学上绝对最小，而要求：
\emph{在既定同伦等价口径与证书基底下，$\mathsf{SemSig}$ 足以唯一锚定同伦类语义（并可回放审计），同时避免携带表层风格噪声。}
\end{remark}

\subsection{投影（压缩）算子 \texorpdfstring{$\pi$}{pi}：从表达到账本可复验签名}
\label{subsec:projection-pi}

\begin{definition}[投影（压缩）算子]
定义一个（部分）投影算子
\[
\pi:\mathsf{Expr}\to \mathsf{SemSig}\sqcup\{\bot\},
\]
其实现是第\ref{sec:multi-scan-convergence}节的多轮扫描 + 第\ref{sec:mdq-granular-audit}节的证书复验与审计落盘：
若收敛并通过复验，则输出语义签名 $s\in\mathsf{SemSig}$；否则输出失败 $\bot$（对应“歧义固定点/空集合/策略拒绝”等）。
\end{definition}

\begin{proposition}[同伦不变性要求（设计约束）]
在目标口径下，投影算子应满足同伦不变性：
\[
x\sim y\ \Longrightarrow\ \pi(x)=\pi(y)
\quad(\text{或至少}\ H_{\mathrm{sem}}(x)=H_{\mathrm{sem}}(y)).
\]
\end{proposition}

\begin{remark}[解释]
该命题不是“自动成立”的数学事实，而是系统应当通过“标准化 + 判别器 + 回归集”去逼近并维护的设计约束：
等价关系 $\sim$ 由体系定义，而 $\pi$ 是该体系对 $\sim$ 的可执行刻画。
\end{remark}

\subsection{展开（表述生成）算子 \texorpdfstring{$G$}{G}：风格约束下的代表元生成}
\label{subsec:expansion-generator}

\begin{definition}[风格空间]
令 $\Sigma$ 为风格约束空间，元素可表示“简洁/正式/口语/学术/客服/法律化/多语言本地化”等生成约束，
并可携带长度、术语表、模板、禁用词等规则。
\end{definition}

\begin{definition}[展开算子（提案器）]
定义展开算子
\[
G:\mathsf{SemSig}\times \Sigma\to \mathsf{Expr},
\]
其实现可以是模板、规则系统或 LLM。关键约束是：$G$ 只负责生成候选表述（提案），
候选是否“仍属于签名所锚定的同伦类”必须由 $\pi$ 与 MDQ 证书体系确定性复验。
\end{definition}

\begin{remark}[“展开无穷多表述”]
对固定签名 $s$ 与风格集合 $\Sigma_0\subseteq \Sigma$，
\[
\{G(s,\sigma):\sigma\in \Sigma_0\}
\]
可视为同一同伦类内的不同代表元（不同表层风格/措辞/格式）。这既可用于合规的风格迁移与鲁棒性测试，
也在结构上覆盖“洗稿式改写”的形式；差别不在数学结构，而在用途合规性与审计归因。
\end{remark}

\subsection{回环一致性（Round-trip 不变性）：\texorpdfstring{$\pi(G(s,\sigma))=s$}{pi(G)=s}}
\label{subsec:roundtrip-invariance}

\begin{definition}[回环一致性（Round-trip）]
称系统满足回环一致性，若对任意可接受签名 $s\in\mathsf{SemSig}$ 与任意风格约束 $\sigma\in\Sigma$，
都有
\[
\pi(G(s,\sigma))=s.
\]
\end{definition}

\begin{remark}[可控生成的本质]
回环一致性把“生成像不像”替换为“投影后是否一致”：不是看表面文字是否相似，而是看投影后的同伦类指纹与参数是否一致。
因此它同时给出：
\begin{itemize}
  \item \textbf{防伪/防洗稿锚点：}同一语义签名的改写难以洗掉指纹；
  \item \textbf{可控展开：}生成出的候选若不能回到同一签名，则自动拒绝或触发修复（PACER 的 Repair）。
\end{itemize}
\end{remark}

\begin{definition}[拒绝与修复（示意）]
给定签名 $s$ 与风格 $\sigma$，生成候选 $x'=G(s,\sigma)$，计算 $s'=\pi(x')$：
\begin{enumerate}
  \item 若 $s'=\bot$，则拒绝并返回失败原因码（来自证书裁决）；
  \item 若 $s'\neq s$，则记录差分事件并触发修复：调整生成约束/补齐缺失参数/升级到更强 witness，再重试；
  \item 若 $s'=s$，则接受并输出，同时落盘审计证据链。
\end{enumerate}
\end{definition}

\subsection{为何等价于 HACA-PACER 最小闭环骨架}
\label{subsec:why-haca-pacer}

\begin{remark}[HACA（分层语义结构）]
MDQ 天然给出分层最小单元（第\ref{subsec:align-haca}小节）。$\pi$ 的实现等价于逐层剥离风格噪声、保留可复验结构骨架，
最终收敛到同伦类语义签名；这就是 HACA 作为“分层语义结构”的可执行版本。
\end{remark}

\begin{remark}[PACER（控制环）]
多次扫描与回环一致性共同构成一个确定性控制过程：
\begin{itemize}
  \item Plan：从审计数据中选出最不稳定/歧义最大的 MDQ 层与类型；
  \item Act：调用（或补齐）对应的证书判别器/标准化器；
  \item Check：用 $\pi$ 与证书验证器复验是否回到固定点/是否满足回环一致性；
  \item Explain/Export：输出证书链、哈希与差分原因码；
  \item Repair：不通过则升级 witness 或修订同伦等价边界/参数口径。
\end{itemize}
这正是 PACER 的“规划--行动--检查--解释--修复”闭环，而且核心裁决不依赖训练。
\end{remark}

\subsection{用途边界：合规风格迁移与不合规洗稿的区分}
\label{subsec:compliance-boundary}

\begin{remark}[结构相同，用途不同]
从纯结构角度，“同一语义签名 $s$ 的多种表述”确实等价于在同伦类内做风格扰动；在形式上与“洗稿”相同。
但工程治理必须区分用途：
\begin{itemize}
  \item 合规用途：风格迁移、可读性改写、模板化输出、多语言本地化、鲁棒性测试、数据增强、对抗样本生成、检索归一化；
  \item 不合规用途：在不保留来源与署名的情况下，用改写逃避查重、规避版权或冒充原创。
\end{itemize}
该体系的反身价值在于：即使有人试图把它当“洗稿器”，也可用 $H_{\mathrm{sem}}$ 与局部指纹集合（第\ref{subsec:layer-local-fingerprint}小节）将其识别出来并形成审计证据。
\end{remark}

\subsection{核心增益（工程化归纳）}
\label{subsec:engineering-gains}

\begin{itemize}
  \item \textbf{跨 LLM 可移植：}LLM 仅作提案器；收敛与签名由 MDQ 证书裁决（第\ref{sec:mdq-portable-layer-homotopy-membership}节）。
  \item \textbf{语义最小化存储与索引：}将海量文本压缩为 $\mathsf{SemSig}$，去重/聚类/统计与检索可直接基于 $H_{\mathrm{sem}}$。
  \item \textbf{生成--复验闭环：}任何生成文本必须通过回环一致性检验，不通过则拒绝或修复（第\ref{subsec:roundtrip-invariance}小节）。
  \item \textbf{可审计、可追责、可回放：}每次差分、证书调用与版本漂移都有证据链（第\ref{sec:mdq-granular-audit}节）。
  \item \textbf{语义不确定的工程化定位：}收敛失败被归结为“边界口径问题”或“基底覆盖问题”（第\ref{subsec:basis-coverage-residual}小节）。
\end{itemize}

\subsection{最硬闭环指标：Round-trip 回归测试规约（建议）}
\label{subsec:roundtrip-regression}

\begin{remark}[Round-trip 不变性作为门槛]
建议将 Round-trip 不变性作为最硬的回归指标：对任意输入 $x$，先压缩 $s=\pi(x)$，再按多种风格生成
$x_i=G(s,\sigma_i)$，要求全部满足 $\pi(x_i)=s$。
该指标长期稳定意味着系统同时具备“可复验压缩收敛”与“可控一致展开”的能力。
\end{remark}

\clearpage

%%%%%%%%%%%%%%%%%%%%%%%%%%%%%%%%%%%%%%%%%%%%%%%%%%%%%%%%%%%%
\section{规范闭环：从最小压缩、不失真展开到语义算子搜索与“意识空间”动力学}
\label{sec:normative-loop-consciousness-dynamics}
%%%%%%%%%%%%%%%%%%%%%%%%%%%%%%%%%%%%%%%%%%%%%%%%%%%%%%%%%%%%

\begin{remark}[核心陈述（可形式化的用途链）]
“规范（语义规范场 + MDQ 证书库 + HACA-PACER 控制环）”的用途可被连成一条可审计链：
\begin{quote}
先把任意风格的表达压到可复验的最小语义签名（最小压缩），
再从该签名在多风格下展开并做回投影一致性检验（不失真），
再在可控的语义算子上做变换与搜索（发现新语义/新算子），
最后把这一切作为一套可审计的“语义状态空间（意识空间）”动力学。
\end{quote}
本章把这条链写成工程可落地的形式系统，并在第\ref{subsec:two-axioms}小节给出两条系统公理（回投影不变性与证书裁决优先）。
\end{remark}

\subsection{核心对象：表达空间、标准化与同伦等价的商空间}
\label{subsec:core-objects-quotient}

\begin{definition}[表达空间与标准化算子]
定义表达空间为
\[
\mathcal{X}:=\{x:\text{任意输入表达}\},
\]
其中 $x$ 可以是自然语言、混合符号串、日志片段、API 入参字符串等。
定义标准化算子
\[
C:\mathcal{X}\to \mathcal{X}_c,
\]
其中 $\mathcal{X}_c$ 是被清洗/规范化后的表达空间。
\end{definition}

\begin{remark}[标准化是“可控地基”】【工程化解释】]
$C$ 是确定性且可复验的步骤：全半角统一、大小写策略、空白/标点归一、数字与单位规范、编码修正等。
它是防伪与可审计的第一道地基，也是后续同伦等价口径的输入标准形。
\end{remark}

\begin{definition}[同伦等价关系]
在 $\mathcal{X}_c$ 上定义同伦等价关系 $\sim$：
\[
x\sim y\quad \Longleftrightarrow\quad
\text{$x$ 与 $y$ 在允许的改写群/规则族下语义等价}.
\]
\end{definition}

\begin{remark}[边界口径]
允许的改写越宽，等价类越大；越窄，等价类越细。
本体系的关键不是“把 $\sim$ 定得一次到位”，而是让 $\sim$ 的维护可以被审计/回归驱动地持续精修（参见第\ref{subsec:fingerprint-boundary}与第\ref{subsec:differential-audit}小节）。
\end{remark}

\begin{definition}[语义签名空间（商空间）]
定义语义签名空间（同伦类集合）为
\[
\mathcal{S}:=\mathcal{X}_c/\sim.
\]
对任意 $x\in\mathcal{X}$，定义类投影（压缩到同伦类）为
\[
\pi_0(x):=[C(x)]_{\sim}\in \mathcal{S}.
\]
\end{definition}

\begin{remark}[与前文记号的对齐]
$\mathcal{S}$ 可与第\ref{subsec:homotopy-class-predicate}小节的“表达同伦类集合”口径对齐；
而第\ref{subsec:projection-pi}小节的 $\pi:\mathsf{Expr}\to\mathsf{SemSig}\sqcup\{\bot\}$ 可视为在 $\pi_0$
  之上补充参数与双层哈希后的\emph{签名级投影}。
\end{remark}

\subsection{MDQ：可外包证书化的离散微分基底}
\label{subsec:mdq-basis-delta}

\begin{definition}[MDQ 基底]
定义 MDQ（微分动力量子）的离散基底集合
\[
\mathcal{D}:=\{\delta_1,\delta_2,\ldots\}.
\]
其中每个 $\delta\in\mathcal{D}$ 是一个可外包证书化的局部算子四元组：
\[
\delta=(\mathrm{span},\ \mathrm{type},\ V_\delta,\ R_\delta).
\]
\end{definition}

\begin{remark}[四元组的工程含义]
\begin{itemize}
  \item $\mathrm{span}$：作用位置/片段（可用 \code{span_digest} 指代）；
  \item $\mathrm{type}$：类型标签（日期、单位、实体、模板片段、逻辑量词等，参见第\ref{subsec:align-haca}小节分层）；
  \item $V_\delta$：验证器（谓词/正则/自动机/小解析器），必须确定性可复验；
  \item $R_\delta$：标准化重写（把片段推到标准形，等价于局部 canonicalization）。
\end{itemize}
这正是“MDQ 粒度审计”的基础：每个 $\delta$ 都可单独验、单独留痕、单独回放（第\ref{sec:mdq-granular-audit}节）。
\end{remark}

\subsection{语义最小压缩：从 \texorpdfstring{$x$}{x} 到 \texorpdfstring{$\pi(x)$}{pi(x)} 的最小可复验签名}
\label{subsec:minimal-compression-semsig}

\begin{definition}[语义最小压缩（签名级投影）]
令 $\mathsf{SemSig}$ 为第\ref{subsec:semsig-space}小节的语义签名空间。
“语义最小压缩”对应签名级投影
\[
\pi:\mathcal{X}\to \mathsf{SemSig}\sqcup\{\bot\},
\]
其实现为“多轮扫描收敛 + MDQ 证书复验 + 审计落盘”（第\ref{sec:multi-scan-convergence}与第\ref{sec:mdq-granular-audit}节）。
\end{definition}

\begin{remark}[最小压缩的目标]
压缩的目标是：把风格/措辞/格式差异压掉，仅保留“同伦类语义锚点 + 必要参数 + 指纹”。
其中 $H_{\mathrm{sem}}$ 提供跨版本稳定的语义锚点，$H_{\mathrm{cert}}$ 提供可回放的证书实现锚点（第\ref{sec:homotopy-fingerprint-hash}节）。
\end{remark}

\subsection{不失真展开：多风格生成与回投影一致性}
\label{subsec:lossless-expansion}

\begin{definition}[风格空间]
定义风格空间为
\[
\Sigma:=\{\sigma_{\text{通俗}},\sigma_{\text{学术}},\sigma_{\text{客服}},\sigma_{\text{法律}},\ldots\},
\]
风格只约束表述方式，不应改变同伦类归属。
\end{definition}

\begin{definition}[展开算子（生成器）]
定义展开算子
\[
G:\mathsf{SemSig}\times \Sigma\to \mathcal{X},
\qquad
x_{\sigma}=G(s,\sigma).
\]
该算子可以由模板或 LLM 实现生成，但其输出必须受 MDQ 证书复验与回投影一致性约束。
\end{definition}

\begin{axiom}[语义守恒（回投影不变性）]
\label{ax:projection-invariance}
对任意 $s\in\mathsf{SemSig}$ 与任意 $\sigma\in\Sigma$，要求
\[
\pi(G(s,\sigma))=s.
\]
\end{axiom}

\begin{remark}[晦涩 $\to$ 通俗的确定性闭环]
给定晦涩输入 $x$，令 $s=\pi(x)$（若 $s=\bot$ 则直接失败），再生成
\[
x_{\text{通俗}}=G(s,\sigma_{\text{通俗}}),
\]
并强制 $\pi(x_{\text{通俗}})=s$。
这保证“先识别语义、再通俗表达”，而不是让 LLM 直接解释（那将引入不可审计漂移）。
\end{remark}

\begin{remark}[多风格一致性作为冗余证据]
取若干风格 $\sigma_1,\ldots,\sigma_k$，令 $x_i=G(s,\sigma_i)$ 并要求 $\pi(x_i)=s$ 对所有 $i$ 成立。
这把“不失真”从单一生成提升为多视角一致性，并自然产生审计证据：哪个风格更易漂移、漂移集中在哪一层 MDQ。
\end{remark}

\subsection{语义变换：在 \texorpdfstring{$\mathcal{S}$}{S} 上操作“语义骨架”而非在文本上乱改}
\label{subsec:semantic-transform-on-quotient}

\begin{definition}[语义变换算子]
定义语义变换算子为
\[
T:\mathcal{S}\to \mathcal{S}.
\]
\end{definition}

\begin{definition}[意义保持型与意义改变型变换]
\begin{enumerate}
  \item 意义保持型（同伦不变）变换满足 $T_{\mathrm{inv}}(s)=s$；
  \item 意义改变型（同伦跃迁）变换满足 $T_{\mathrm{chg}}(s)=s'$ 且 $s'\neq s$。
\end{enumerate}
\end{definition}

\begin{remark}[审计要求]
风格迁移/同义改写应落入意义保持型；而“抽象化/具体化、加入条件/删除条件、量词强度改变”等应显式表现为意义改变型跃迁。
系统必须能说清楚“语义变了”，而不是悄悄变；否则审计与治理不成立。
\end{remark}

\subsection{语义规范场：联络写法、离散协变微分与上同调障碍的工程化记账}
\label{subsec:semantic-gauge-field-connection}

\begin{definition}[语义域覆盖与联络（示意）]
令 $\{U_\alpha\}$ 表示语义域的覆盖（不同语境/子领域/模板协议）。
在每个 $U_\alpha$ 上存在局部表示与局部证书约束。定义一个联络形式写法
\[
D:=d+A,
\]
其中 $d$ 表示基础差分/变分，$A$ 表示语义规范场的纠偏项（由“标准最小语义函子族/证书库”实现的可控算子族）。
\end{definition}

\begin{definition}[MDQ 作为离散协变微分（离散边缘算子）]
考虑一条离散语义轨线 $s_0\to s_1\to\cdots\to s_T$，令
\[
s_{t+1}=s_t\oplus \delta_t,
\qquad
\delta_t\in\mathcal{D}.
\]
其中 $\oplus$ 表示“通过一个可证书化的 MDQ 量子对语义骨架做一步局部更新”。
\end{definition}

\begin{remark}[PACER 的可修复性落点]
当局部规则/不同层级约束无法同时满足时，可记账为障碍（残差）出现：
\[
\mathrm{Ob}(x)\neq 0.
\]
PACER 的职责是：选择或引入合适的 $\delta$（或扩充 $A$ 的最小函子族）把障碍压到 $0$；
若在受控资源与安全子集约束下无法消除，则返回失败（\code{ok:false}）并把原因码写入审计事件流。
\end{remark}

\subsection{GRL 路径积分：在语义轨线上探索与固化“新同伦类/新证书”}
\label{subsec:grl-path-integral}

\begin{definition}[语义轨线集合]
给定初始语义 $s_0\in\mathcal{S}$，定义从 $s_0$ 出发的轨线集合
\[
\Gamma(s_0):=\{\gamma: s_0\to s_T\},
\]
其中每条轨线 $\gamma$ 可写为 MDQ 更新序列 $\gamma=(\delta_0,\ldots,\delta_{T-1})$。
\end{definition}

\begin{definition}[作用量（治理标量）]
对轨线 $\gamma$ 定义作用量
\[
\mathcal{A}(\gamma)=\sum_{t=0}^{T-1}\Bigl(
\lambda_1\,\mathrm{Amb}(s_t)
\;+\;
\lambda_2\,\mathrm{Cost}(\delta_t)
\;+\;
\lambda_3\,\mathrm{Risk}(s_t)
\;-\;
\lambda_4\,\mathrm{Novel}(s_t)
\Bigr),
\]
其中 $\mathrm{Amb}$ 表示歧义度，$\mathrm{Cost}$ 表示验证/表达成本，$\mathrm{Risk}$ 表示误判或不可审计风险，
$\mathrm{Novel}$ 表示新语义/新结构增益，$\lambda_i\ge 0$ 为治理参数。
\end{definition}

\begin{definition}[路径积分分布与终点权重（示意）]
定义轨线的加权分布为
\[
\mathbb{P}(\gamma)\propto \exp\bigl(-\mathcal{A}(\gamma)\bigr),
\]
并定义终点的聚合权重为
\[
p(s_T)=\sum_{\gamma\in\Gamma(s_0),\ \gamma(T)=s_T}\exp\bigl(-\mathcal{A}(\gamma)\bigr).
\]
\end{definition}

\begin{remark}[“发现新语义”的工程解释]
若某些终点 $s_T$ 不属于现有同伦类库，但在低作用量路径上高频出现（$p(s_T)$ 较大），
则可将其固化为新的同伦类/新的谓词/新的 MDQ 证书（即“新语义算子”）。
该过程的关键仍是：所有步进与固化必须可证书化、可回放、可审计。
\end{remark}

\subsection{“意识空间”等价：把 \texorpdfstring{$(\mathcal{S},\mathcal{O})$}{(S,O)} 视为可审计的状态空间动力学}
\label{subsec:consciousness-space-modeling}

\begin{definition}[意识空间（语义状态空间）]
定义（建模意义上的）意识空间为二元组
\[
\mathcal{C}:=(\mathcal{S},\ \mathcal{O}),
\]
其中 $\mathcal{S}$ 为语义状态集合（同伦类），$\mathcal{O}$ 为在其上作用的算子集合（MDQ 步进、风格展开、语义变换、修复算子、搜索算子等）。
\end{definition}

\begin{definition}[语义不确定性状态（分布）]
令
\[
\Delta(\mathcal{S}):=\Bigl\{\psi:\mathcal{S}\to [0,1]\ :\ \sum_{s\in\mathcal{S}}\psi(s)=1\Bigr\}
\]
为 $\mathcal{S}$ 上的概率单纯形（离散口径）。
取 $\psi\in\Delta(\mathcal{S})$ 表示“当前对语义归属的不确定性分布”。
\end{definition}

\begin{definition}[演化（控制环驱动）]
定义由 PACER/搜索/审计共同诱导的演化算子
\[
\psi_{t+1}=\mathcal{T}(\psi_t),
\]
其中 $\mathcal{T}$ 由“提案（LLM）+ 复验（MDQ）+ 修复（证书增补/边界收紧）+ 记录（审计）+ 搜索（路径积分）”组合实现。
\end{definition}

\begin{remark}[建模等价而非经验断言]
以上“意识空间”仅表达一种\emph{建模等价}：把语义签名空间与算子代数视为一个可审计的状态空间动力学系统，
用以组织“从多解到单解”的收敛过程；并不对人脑意识作经验性断言。
\end{remark}

\subsection{用途对齐：把工程需求投影到上述结构}
\label{subsec:use-case-projection}

\begin{remark}[用途 $\leftrightarrow$ 结构对应]
\begin{itemize}
  \item 晦涩识别后通俗展开且不失真：$s=\pi(x)$，$x_{\text{通俗}}=G(s,\sigma_{\text{通俗}})$，并强制 $\pi(x_{\text{通俗}})=s$。
  \item 多风格展开确保不失真：对多个 $\sigma_i$ 展开并要求全部回投影一致；漂移留痕并触发修复。
  \item 语义变换：在 $\mathcal{S}$ 上施加 $T$，并明确标注 $T_{\mathrm{inv}}$ 或 $T_{\mathrm{chg}}$，使语义改变可审计。
  \item 基于语义规范场与路径积分发现新语义：用 $\mathcal{D}$ 步进定义可控轨线，用 $\exp(-\mathcal{A}(\gamma))$ 选择低成本高新颖路径，
  将稳定终点固化为新同伦类与新证书。
  \item “意识空间”动力学：$(\mathcal{S},\mathcal{O})$ 上的 $\psi$ 收敛，等价于“从多解到单解”的可审计过程。
  \item 洗稿（同伦类内风格扰动）的双向能力：既能生成（$G$），也能用 $H_{\mathrm{sem}}$ 与证书链识别并审计（第\ref{sec:homotopy-fingerprint-hash}节）。
\end{itemize}
\end{remark}

\subsection{两条系统公理：回投影不变性与证书裁决优先}
\label{subsec:two-axioms}

\begin{axiom}[证书裁决优先（LLM 只做提案）]
\label{ax:certificate-decision-priority}
对任意 $x\in\mathcal{X}$ 与任意 $s\in\mathcal{S}$，定义接受关系 $\mathrm{Accept}(x\in s)$。
要求存在与 $s$ 绑定的确定性标准谓词/判别器 $V_s$（证书化、可复验），使得
\[
\mathrm{Accept}(x\in s)\ \Longleftrightarrow\ V_s\bigl(C(x)\bigr)=1.
\]
\end{axiom}

\begin{remark}[两公理的含义]
公理~\ref{ax:projection-invariance} 保证语义守恒（不失真展开）；公理~\ref{ax:certificate-decision-priority} 保证裁决以证书为准（LLM 可替换）。
二者一旦成立，则通俗展开、多风格一致性、语义变换可审计、路径积分发现新语义、以及“意识空间”动力学口径，都可视为同一机制的不同投影。
\end{remark}

\clearpage

%%%%%%%%%%%%%%%%%%%%%%%%%%%%%%%%%%%%%%%%%%%%%%%%%%%%%%%%%%%%
\section{倒逼机制：把 LLM 提案层变成可度量、可筛选、可淘汰的组件}
\label{sec:selection-pressure-llm-proposal-layer}
%%%%%%%%%%%%%%%%%%%%%%%%%%%%%%%%%%%%%%%%%%%%%%%%%%%%%%%%%%%%

\begin{remark}[核心陈述（倒逼的因果链）]
一旦系统把“是否属于某同伦类”的裁决权从“LLM 说了算”改为“证书谓词 $V_s$ 说了算”，
并把误收/漏收、耗时与成本写入不可变审计日志，那么 LLM 的行为就会被自然倒逼为一个可度量、可筛选、可淘汰的提案层：
不够快、不够稳或不贴近共识标准的模型/提示/路由组合，会因为更高的拒绝率、更频繁的回退与更差的指标而被系统降权或替换。
\end{remark}

本章把上述倒逼写成可执行的符号系统：以证书谓词裁决为“裁判”，以审计指标为“成绩”，以门槛与路由为“选择算子”。
关键点在于：所谓“共识标准”不是 LLM 的共识，而是“标准化器 $C$ + 证书谓词库 $\{V_s\}$ + 回归集”的共识（见第\ref{subsec:two-axioms}小节）。

\subsection{单次请求的记号：裁决、提案与确定性复验}
\label{subsec:single-request-notation}

对任意输入 $x\in\mathcal{X}$，记其同伦类裁决为
\[
s^\ast:=\pi_0(x)\in\mathcal{S},
\]
其中 $\pi_0(x)=[C(x)]_{\sim}$ 见第\ref{subsec:core-objects-quotient}小节。

令 $m$ 表示某个“提案器配置”（模型版本、prompt 版本、路由策略或其组合）。定义其提案输出为
\[
(\hat{s}_m(x),\hat{y}_m(x)):=\mathrm{Prop}_m(x),
\qquad
\hat{s}_m(x)\in\mathcal{S},\ \hat{y}_m(x)\in\{0,1\}.
\]
其中 $\hat{s}_m(x)$ 是其建议的同伦类（或候选集合中的首选），$\hat{y}_m(x)$ 是其对“属于/不属于”的布尔判断。

确定性复验由证书谓词给出：
\[
y^\ast_m(x):=V_{\hat{s}_m(x)}\bigl(C(x)\bigr)\in\{0,1\},
\]
其中 $V_s$ 满足公理~\ref{ax:certificate-decision-priority}（证书裁决优先）。

\begin{remark}[裁决与提案的角色分离]
$\hat{s}_m,\hat{y}_m$ 仅是提案；$y^\ast_m$ 是确定性裁决结果。
任何“质量好坏”都应以 $y^\ast_m$ 为准进行度量与追责，而不是以 LLM 的口头解释为准。
\end{remark}

\subsection{可机械判定的偏差：误收/漏收与成本记账}
\label{subsec:fp-fn-cost}

\begin{definition}[误收与漏收（指示变量）]
定义两类关键偏差为：
\[
\mathrm{FP}_m(x):=\mathbf{1}\bigl[\hat{y}_m(x)=1\ \wedge\ y^\ast_m(x)=0\bigr],
\qquad
\mathrm{FN}_m(x):=\mathbf{1}\bigl[\hat{y}_m(x)=0\ \wedge\ y^\ast_m(x)=1\bigr].
\]
\end{definition}

\begin{remark}[审计落盘字段]
只要把 $(\mathrm{FP}_m(x),\mathrm{FN}_m(x))$ 与调用成本、耗时、失败原因码一并写入审计日志（第\ref{sec:mdq-granular-audit}节），
就获得了可回放、可对比、可聚合统计的“适应度信号”，从而使倒逼成为一种可执行的制度而非口头约束。
\end{remark}

\begin{definition}[耗时、成本与回退次数（示意）]
记 $t_m(x)$ 为端到端延迟（包含必要的多轮扫描与回退），$c_m(x)$ 为成本（调用次数、token、额外验证开销等），
并记 $r_m(x)$ 为“回退/重扫”次数（例如提案失败触发的重提案次数）。
\end{definition}

\begin{remark}[与“注意力/专注”记账口径的对齐]
将 $t_m,c_m,r_m$ 作为路由代价项，本质上是在工程层面对“计算注意力/专注预算”做可观测化与可治理化；
其思想与 Vol3 中用占用测度刻画“语义注意力/时间/频率”分配的口径一致\cite{GaoZhengAppVol3}。
\end{remark}

\subsection{倒逼的抓手：在真实分布上的损失函数与选择规则}
\label{subsec:loss-and-selection}

\begin{definition}[请求分布与损失函数]
令 $D$ 为请求分布。对任意提案器配置 $m$，定义其总体损失为
\[
\mathcal{L}(m)=
\lambda_{\mathrm{FP}}\ \mathbb{E}_{x\sim D}[\mathrm{FP}_m(x)]
\;+\;
\lambda_{\mathrm{FN}}\ \mathbb{E}_{x\sim D}[\mathrm{FN}_m(x)]
\;+\;
\lambda_t\ \mathbb{E}_{x\sim D}[t_m(x)]
\;+\;
\lambda_c\ \mathbb{E}_{x\sim D}[c_m(x)]
\;+\;
\lambda_r\ \mathbb{E}_{x\sim D}[r_m(x)],
\]
其中 $\lambda_{\mathrm{FP}},\lambda_{\mathrm{FN}},\lambda_t,\lambda_c,\lambda_r\ge 0$ 为治理参数。
\end{definition}

\begin{remark}[权重的治理含义]
安全/合规场景通常取 $\lambda_{\mathrm{FP}}>\lambda_{\mathrm{FN}}$（宁可漏判也不误判）；
召回优先场景可相反。无论权重如何取，倒逼的核心是：把“好坏”落实为可复验的统计量，而非口头偏好。
\end{remark}

\begin{definition}[选择规则（最朴素版本）]
若候选提案器集合为 $\mathcal{M}$，则可取
\[
m^\ast=\arg\min_{m\in\mathcal{M}}\ \mathcal{L}(m).
\]
也可在门槛约束下做受限选择，例如要求 $\mathbb{E}[\mathrm{FP}_m]\le \alpha$、$\mathbb{E}[t_m]\le T_{\max}$ 等。
\end{definition}

\begin{proposition}[倒逼成立（度量 $\Rightarrow$ 淘汰）]
在固定证书标准 $\{V_s\}$、固定标准化器 $C$ 与固定审计口径下，
若系统按上述损失函数与门槛规则进行路由选择，则任何导致更高误收/漏收、更高回退次数或更高成本/延迟的提案器配置，
都会在统计意义上被降权或淘汰。
\end{proposition}

\subsection{为何会更快：候选集合上限、早停缓存与回退罚分}
\label{subsec:faster-pressure}

\begin{definition}[候选集合大小约束（示意）]
允许提案器输出候选集合 $\widehat{S}_m(x)\subseteq\mathcal{S}$，并施加硬约束
\[
|\widehat{S}_m(x)|\le K.
\]
若超限，则直接拒绝或计入罚分（体现为 $c_m(x)$ 或 $r_m(x)$ 上升）。
\end{definition}

\begin{remark}[为何能加速]
候选集合越大，验证成本近似线性上升；强制 $|\widehat{S}|\le K$ 会倒逼提案器做更聚焦的提案，
否则它会在 $c_m,t_m,r_m$ 指标上持续吃亏。
\end{remark}

\begin{remark}[早停与缓存（系统能力）]
若命中语义指纹缓存（第\ref{sec:homotopy-fingerprint-hash}节），可采用早停策略：
\begin{quote}
若 $H_{\mathrm{sem}}(x)$ 命中缓存，则跳过提案器，直接复用已审计的签名/证书链。
\end{quote}
这会把“快”从模型能力提升为系统能力：系统运行越久、缓存越大，整体越快。
\end{remark}

\subsection{为何会更稳：同义变体一致性与提案漂移率}
\label{subsec:stability-pressure}

\begin{definition}[同义变体集合与漂移率（示意）]
对同一语义输入 $x$，令 $\mathcal{V}(x)=\{x_1,\ldots,x_n\}$ 为其同义改写/噪声扰动集合（第\ref{subsec:metamorphic-testing}小节）。
定义提案漂移率为
\[
\mathrm{Drift}(m):=\mathbb{E}_{x\sim D}\Bigl[\mathbf{1}\bigl(\exists i,j,\ \hat{s}_m(x_i)\neq \hat{s}_m(x_j)\bigr)\Bigr]
  .
\]
\end{definition}

\begin{remark}[稳定的含义]
“稳”在本体系中等价于：在同义改写、风格扰动与噪声扰动下，提案与裁决的一致性更高；
漂移率高的提案器会导致更多回退与不一致审计事件，从而在路由上被降权。
\end{remark}

\subsection{为何会更符合共识标准：裁判是 \texorpdfstring{$(C,\{V_s\},\text{回归集})$}{(C,Vs,regression)}}
\label{subsec:consensus-pressure}

\begin{remark}[“共识标准”被固化为证书体系]
由于接受关系由公理~\ref{ax:certificate-decision-priority} 固化为 $V_s(C(x))$，
提案器只有两条路：
\begin{itemize}
  \item 学会“按标准提案”：更准确预测哪些 $V_s$ 会通过，并给出更贴近 $C$ 的调用顺序；
  \item 持续制造 $\mathrm{FP}/\mathrm{FN}$ 与回退事件，在审计中不断暴露并被路由系统替换。
\end{itemize}
因此，“更符合共识”不是道德约束，而是带裁判的竞技场：裁判是证书谓词库，成绩是审计指标。
\end{remark}

\subsection{边界条件：共识标准的治理与双层指纹的作用}
\label{subsec:governance-boundary}

\begin{remark}[避免“把错误标准执行得很一致”】【治理要求】]
倒逼成立的前提是：共识标准本身具备治理机制，否则可能出现“把错误标准执行得很一致”。
工程上通常用双层指纹保证可解释性：
\begin{itemize}
  \item $H_{\mathrm{sem}}$ 尽量跨版本稳定，用于长期索引、防洗稿与聚类；
  \item $H_{\mathrm{cert}}$ 明确绑定谓词/验证器/标准化器版本，用于回放复验与漂移归因。
\end{itemize}
这样即使标准更新，也能区分“标准变了”与“提案器变差了”，并将变化落到可审计证据链上。
\end{remark}

\subsection{可上线的最小 KPI 门槛（建议）}
\label{subsec:min-kpi-gates}

\begin{remark}[从理念到制度]
若要把倒逼机制“坐实”为可上线制度，可设置一组最小 KPI 与门槛（仅示意）：
\begin{itemize}
  \item $\mathbb{E}[\mathrm{FP}_m]\le \alpha_{\mathrm{FP}}$，$\mathbb{E}[\mathrm{FN}_m]\le \alpha_{\mathrm{FN}}$；
  \item $\mathbb{E}[t_m]\le T_{\max}$，$\mathbb{E}[c_m]\le C_{\max}$；
  \item $|\widehat{S}_m(x)|\le K$（硬约束），$\mathbb{E}[r_m]\le R_{\max}$；
  \item $\mathrm{Drift}(m)\le \delta_{\max}$（稳健性门槛）。
\end{itemize}
门槛触发即降权或淘汰，并由差分审计输出解释（第\ref{subsec:differential-audit}小节）。
\end{remark}

\clearpage

%%%%%%%%%%%%%%%%%%%%%%%%%%%%%%%%%%%%%%%%%%%%%%%%%%%%%%%%%%%%
\section{标准先行、多实现：OS+JVM 范式下的 MDQ 可移植语义层}
\label{sec:os-jvm-paradigm}
%%%%%%%%%%%%%%%%%%%%%%%%%%%%%%%%%%%%%%%%%%%%%%%%%%%%%%%%%%%%

\begin{remark}[核心陈述（结构骨架一致）]
将“MDQ 证书外包 + HACA-PACER 收敛裁决 + 同伦类哈希”的体系类比为 \textsf{OS+JVM}（或“标准先行、厂商多实现”）范式，
关键不在于“像不像”，而在于共享同一条工程骨架：
\begin{itemize}
  \item 标准定义一个稳定的中间层（IR/指令集/可验证格式）；
  \item 多厂商实现可自由优化上层（编译器/JIT/调度器），但必须通过标准验证与一致性测试；
  \item 最终保证可移植与可审计：实现可变，语义锚点不变。
\end{itemize}
本章给出一个一一对应的形式化映射，并把一致性条件写成可执行合同。
\end{remark}

\subsection{稳定中间层：MDQ 调用序列作为可验证 IR}
\label{subsec:mdq-as-ir}

沿用第\ref{subsec:mdq-basis-delta}小节：$\mathcal{D}$ 为 MDQ 基底集合。
记 $\mathcal{D}^{\ast}$ 为 MDQ 序列集合（Kleene 星），并把一条序列写为
\[
\gamma=(\delta_0,\ldots,\delta_{T-1})\in\mathcal{D}^{\ast}.
\]

\begin{definition}[MDQ 调用序列（IR）的可验证性]
称一条 MDQ 序列 $\gamma$ \emph{可验证}，若其每个 $\delta_t$ 的证书谓词与标准化重写均可被确定性复验，
并且全序列在策略快照（白名单、复杂度阈值等）下通过一致性检查。
\end{definition}

\begin{remark}[类比]
这对应 \textsf{JVM} 中“字节码格式 + bytecode verifier”：
生成字节码的人可以很多，但进入运行时之前必须能被验证与执行，且语义明确。
在本体系中，LLM/路由/提示都可以生成不同的 MDQ 调用序列，但进入裁决层之前必须通过证书复验与审计约束（参见 \cite{jvms}）。
\end{remark}

\subsection{多实现的“编译产物”与标准执行器：\texorpdfstring{$\mathrm{Comp}_m$}{Comp_m} 与 \texorpdfstring{$\mathrm{Eval}$}{Eval}}
\label{subsec:comp-eval-contract}

令表达空间为 $\mathcal{X}$，标准化为 $C:\mathcal{X}\to\mathcal{X}_c$，同伦等价为 $\sim$，
并令语义签名空间为 $\mathcal{S}=\mathcal{X}_c/\sim$（第\ref{subsec:core-objects-quotient}小节）。
记 $\pi_0(x)=[C(x)]_{\sim}\in\mathcal{S}$ 为同伦类投影。

\begin{definition}[多厂商编译器：从表达生成 MDQ 序列]
对任意实现/厂商/模型配置 $m$（不同 LLM、不同 prompt、不同推理引擎或其组合），定义其“编译器”为
\[
\mathrm{Comp}_m:\mathcal{X}\to \mathcal{D}^{\ast},
\]
其输出 $\gamma=\mathrm{Comp}_m(x)$ 可理解为对输入 $x$ 的一条“MDQ 调用提案序列”（类比字节码/IR）。
\end{definition}

\begin{definition}[标准执行与裁决：确定性评估器]
定义确定性评估器（标准执行器）为一个（部分）函数
\[
\mathrm{Eval}:\mathcal{D}^{\ast}\to \mathcal{S}\sqcup\{\bot\},
\]
其作用是：对输入 MDQ 序列做白名单/一致性检查，并按证书谓词与标准化重写确定性执行；
若序列可验证并收敛，则输出语义签名 $s\in\mathcal{S}$；否则输出失败 $\bot$ 并给出原因码（在实现层落盘审计）。
\end{definition}

\begin{proposition}[可移植性合同：\texorpdfstring{$\mathrm{Eval}\circ \mathrm{Comp}_m=\pi_0$}{Eval∘Comp=pi0}]
\label{prop:portability-contract}
对任意 $m$ 与任意 $x\in\mathcal{X}$，若 $\mathrm{Eval}(\mathrm{Comp}_m(x))\neq \bot$，则必须满足
\[
\mathrm{Eval}\bigl(\mathrm{Comp}_m(x)\bigr)=\pi_0(x).
\]
\end{proposition}

\begin{remark}[文字解释]
不管 $m$ 如何实现、如何加速、如何优化，只要它产出的 MDQ 调用能被标准执行器复验并收敛，
最终就必须回到同一个语义签名 $\pi_0(x)$。这就是“可移植”：上层实现可变，裁决结果（语义锚点）不变。
若某实现频繁导致 $\bot$（不可验证/不收敛/超策略约束），则会在第\ref{sec:selection-pressure-llm-proposal-layer}节的倒逼机制下被降权或淘汰。
\end{remark}

\subsection{路径不同、同伦类相同：多实现差异的工程版同伦}
\label{subsec:path-homotopy-osjvm}

\begin{remark}[工程版含义]
不同实现可能给出不同的 MDQ 路径（$\gamma_1\neq \gamma_2$），但在标准执行与复验下收敛到同一终点：
\[
\mathrm{Eval}(\gamma_1)=\mathrm{Eval}(\gamma_2)=s.
\]
这可被视为“路径不同、同伦类相同”的工程版含义：差异存在于提案路径与优化策略，
而语义等价由标准裁决层给出。
\end{remark}

\begin{remark}[类比]
这类似于：不同编译器/JIT/调度器可以采取不同优化与执行路径，但只要通过验证与一致性测试，
它们就必须实现同一个被标准承诺的语义结果。
\end{remark}

\subsection{OS+JVM 对应表：证书库、语义锚点、PACER 调度与审计回放}
\label{subsec:osjvm-mapping}

\begin{remark}[部件级对应（要点）]
\begin{enumerate}
  \item \textbf{MDQ 证书/谓词库 $\approx$ 字节码规范 + Verifier 规则.}
  关键不是“谁来生成”，而是“产物必须可验证、可执行、语义明确”。
  \item \textbf{语义签名 $\mathcal{S}$ 与 $H_{\mathrm{sem}}$ $\approx$ 平台无关语义锚点.}
  任意风格改写只要同伦等价，就应落到同一 $H_{\mathrm{sem}}$（第\ref{sec:homotopy-fingerprint-hash}节）。
  \item \textbf{PACER 调度/回退 $\approx$ OS 调度器 + JVM 运行时（解释/JIT/缓存/早停）.}
  在满足合同（命题~\ref{prop:portability-contract}）的前提下追求提速与降本。
  \item \textbf{审计留痕与 $H_{\mathrm{cert}}$ $\approx$ 运行时日志/堆栈/崩溃转储 + 可回放一致性.}
  可定位到“哪个 MDQ、哪个证书版本、哪个验证器、哪次提案导致分歧”，并支持事后重放（第\ref{subsec:replay-audit}小节）。
\end{enumerate}
\end{remark}

\subsection{标准化基础设施：Spec、参考验证器与一致性测试套件（TCK 类比）}
\label{subsec:spec-verifier-tck}

\begin{remark}[三件套]
要获得类似 \textsf{JVM} 生态的“多实现可互换”效果，通常需要三类标准化基础设施（在本体系中均有直接对应）：
\begin{enumerate}
  \item \textbf{规范文档（Spec）.}
  为每个同伦类/MDQ 类型写清边界、允许的标准化、失败条件（何时必须返回 \code{ok:false}）。
  \item \textbf{参考验证器（Reference Verifier）.}
  确定性执行器/校验器是“唯一裁判”，并随版本发布其白名单与回归策略快照。
  \item \textbf{一致性测试集（Conformance Suite / TCK）.}
  用固定回归集与指标门槛比较不同实现的 $\mathrm{FP}/\mathrm{FN}$、漂移率、延迟与成本，并规定最低合格线。
\end{enumerate}
（TCK 的工程口径可参见 \cite{java-tck}。）
\end{remark}

\begin{remark}[与倒逼效应的连接]
当一致性测试套件成为“准入门槛”，并与路由权重/预算约束绑定时，
系统就会自然产生倒逼效应：实现为了通过测试并获得更高权重，必须输出更贴近标准、更精简、更高命中率的 MDQ 提案（第\ref{sec:selection-pressure-llm-proposal-layer}节）。
\end{remark}

\subsection{工程口号（摘要）}
\label{subsec:osjvm-slogan}

\begin{remark}[口号]
\textsf{JVM}: \emph{Write once, run anywhere.}\\
本体系更像：\emph{Write any style, verify to one semantics} —— 任意风格输入都能被压到同一可复验语义锚点；
任意厂商 LLM 都可以参与提案，但必须接受同一裁决与审计。
\end{remark}

\clearpage

%%%%%%%%%%%%%%%%%%%%%%%%%%%%%%%%%%%%%%%%%%%%%%%%%%%%%%%%%%%%
\section{微型模型生态：围绕 MDQ/HACA-PACER 的专用提案器、策略器与渲染器}
\label{sec:micro-model-ecosystem}
%%%%%%%%%%%%%%%%%%%%%%%%%%%%%%%%%%%%%%%%%%%%%%%%%%%%%%%%%%%%

\begin{remark}[核心判断（自然产物）]
在“标准化语义中间层（MDQ / HACA-PACER）+ 证书裁决（Verifier）+ 一致性测试（Conformance Suite）”的结构下，
出现专门迎合该体系的微型 LLM（或小型 Transformer/策略网络）是高度可能的，且几乎是自然产物。
原因在于：一旦语义系统被压成可复验的中间表示（同伦类签名 / MDQ 调用序列），
对提案器的要求从“生成整段正确文本”收窄为“提出正确的 MDQ 调用与候选同伦类”。
这更接近编译器前端/路由器/标注器，天然适配更小、更专用的模型。
\end{remark}

\subsection{为何会出现“迎合 MDQ/HACA-PACER”的微型模型}
\label{subsec:why-micro-models}

在第\ref{sec:os-jvm-paradigm}节中，我们已将多实现结构写成
\[
\mathrm{Comp}_m:\mathcal{X}\to \mathcal{D}^{\ast},
\qquad
\mathrm{Eval}:\mathcal{D}^{\ast}\to \mathcal{S}\sqcup\{\bot\},
\]
并以命题~\ref{prop:portability-contract} 给出“可验证时必须回到同一语义签名”的合同。
在第\ref{sec:selection-pressure-llm-proposal-layer}节中，我们进一步把提案器配置 $m$ 的好坏度量为
$\mathrm{FP}/\mathrm{FN}$、延迟、成本、回退次数与漂移率等可统计指标，并通过路由/门槛形成选择压力。

\begin{remark}[输出空间收窄 $\Rightarrow$ 小模型优势]
当竞争目标被固化为：
\begin{itemize}
  \item 更小的候选集合上限 $K$（更少无用验证）；
  \item 更低的 $\mathrm{FP}/\mathrm{FN}$（更少误收/漏收）；
  \item 更少回退（更少重扫）；
  \item 更低延迟、更低成本；
\end{itemize}
且输出空间主要由 \code{class_id}/\code{span}/\code{mdq_type}/调用顺序等标准化结构构成时，
大模型擅长的长篇生成能力反而不是瓶颈；小模型更容易在固定预算内做出高性价比的“编译式提案”。
\end{remark}

\subsection{微型模型会长什么样：四类专用形态}
\label{subsec:micro-model-types}

\begin{definition}[A 类：MDQ 路由器（Router）]
MDQ 路由器的目标是把“全库扫描”降为“少量候选复验”。可形式化为
\[
\mathsf{Route}:\mathcal{X}\to \mathcal{P}_{\le K}(\mathcal{S})\times [0,1]^K,
\]
输出 top-$K$ 候选语义签名（或同伦类）及其置信度排序，从而优先触发高收益的证书复验。
\end{definition}

\begin{definition}[B 类：MDQ 片段标注器（Span Tagger）]
MDQ 片段标注器在字符/token 级标出 $\mathrm{span}$ 并给出 $\mathrm{type}$：
\[
\mathsf{Tag}:\mathcal{X}\to (\mathsf{Span}\times \mathsf{Type})^{\ast}.
\]
其价值在于减少 PACER 的探索空间，使“修复”更像确定性流水线而不是盲搜。
\end{definition}

\begin{definition}[C 类：PACER 策略模型（Policy / Scheduler）]
PACER 策略器根据当前状态（哪些 MDQ 已通过、哪里失败、歧义度/成本指标等）选择下一步动作：
\[
\mathsf{Pol}:\mathsf{State}\to \mathsf{Action},
\]
例如“下一步调用哪个 MDQ/是否回退/是否升级 witness 档位”。该模型甚至不必是 LLM；
小型策略网络或规则系统即可，但用小 Transformer 往往更能覆盖多样表达。
\end{definition}

\begin{definition}[D 类：风格展开器（Style Realizer）]
在语义签名 $s$ 已确定的前提下，风格展开器实现
\[
G(s,\sigma)\in\mathcal{X},
\]
并通过公理~\ref{ax:projection-invariance} 的回投影约束（$\pi(G(s,\sigma))=s$）形成闭环。
由于失败可被复验器拦下并触发修复，该展开器可做得更小：它不需要“理解世界”，只需“在固定语义下做风格渲染”。
\end{definition}

\begin{remark}[“迎合”的含义]
所谓“迎合”并非迎合自由对话能力，而是迎合标准裁决体系：只要它能在审计指标上更好地逼近
\[
\mathrm{Eval}(\mathrm{Comp}_m(x))\approx \pi_0(x),
\]
并在 $K$、$\mathrm{FP}/\mathrm{FN}$、回退率与成本/延迟上更优，它就能在路由层获得更高权重。
这与第\ref{sec:os-jvm-paradigm}节的 OS+JVM 生态类比一致：很多厂商会去做更好的编译器/JIT/工具链，而不是“最强语言”。
\end{remark}

\subsection{数据来源：审计留痕与证书回放天然形成监督信号}
\label{subsec:micro-model-data}

本体系已将两类错误（$\mathrm{FP}/\mathrm{FN}$）与失败原因码机械化并写入审计事件流（第\ref{subsec:fp-fn-cost}小节）。
因此每次请求都会产生可监督样本：
\begin{itemize}
  \item 输入：原始表达 $x$ 或标准化表达 $C(x)$（可配合脱敏/摘要策略）；
  \item 标签：通过复验的同伦类（$s^\ast=\pi_0(x)$ 或其 \code{class_id}）、
  以及可验证的 MDQ 序列/片段（$\gamma=\mathrm{Comp}_m(x)$ 在通过时的版本化记录）与失败原因码。
\end{itemize}

\begin{remark}[运行时蒸馏的自然性]
即便早期坚持“不训练也能跑”，系统一旦上线，日志就会自动积累训练集；
随后围绕 Router/Tagger/Policy/Realizer 的蒸馏与小模型训练将变得非常自然。
关键在于：监督标签由证书裁决层给出，而非由 LLM 自述给出，因此更贴近共识标准。
\end{remark}

\subsection{需要提前防的一点：防止“为测试而优化”的伪稳健}
\label{subsec:anti-gaming}

\begin{remark}[风险：刷分与过拟合]
若出现微型模型生态，常见风险是“为测试而优化”：对一致性测试集或某类输入分布过拟合，
导致表面通过但对真实扰动脆弱。此时第\ref{subsec:metamorphic-testing}小节的变形测试从“建议”变成刚需：
\begin{itemize}
  \item 对同一语义的变体集合 $\{x_i\}$，要求收敛一致（$\pi(x_i)$ 相同）；
  \item 对对抗扰动与噪声扰动，要求 $\mathrm{FP}/\mathrm{FN}$ 不显著恶化；
  \item 对新出现表达形态，允许返回失败并触发补证书，而不是硬猜（保持 $\bot$ 的语义边界清晰）。
\end{itemize}
通过动态回归集、对抗样本与多风格 round-trip（第\ref{subsec:roundtrip-regression}小节），可抑制刷分并维持真实稳健性。
\end{remark}

\subsection{总体路线意义：标准先行生态与多小模型协同竞争}
\label{subsec:ecosystem-meaning}

\begin{remark}[对 Vol3 论证的反向强化]
微型模型生态反而会强化 Vol3 的论证：Vol3 给出抽象严谨的语义几何与算子体系；
工程实现层则自然产生“专用前端/专用策略/专用渲染”的分工生态。
系统不需要一个万能模型，而可以让多个小模型围绕同一套证书裁决与语义签名标准协同竞争，
这正是“标准先行、多实现一致”的长期形态（第\ref{sec:os-jvm-paradigm}节）。
\end{remark}

\subsection{把生态导向正确方向：三条强约束（建议）}
\label{subsec:steering-constraints}

\begin{remark}[三条强约束]
若希望生态演进不破坏语义收敛与可审计性，最有效的做法不是指定某个模型，而是把三样东西做成强约束：
\begin{enumerate}
  \item 证书谓词与验证器的参考实现（裁判不可替换）；
  \item 回投影不变性与变形测试构成的 conformance suite（测试不可被刷分）；
  \item 审计指标（$\mathrm{FP}/\mathrm{FN}$、延迟、回退率、覆盖率）作为路由权重依据（选择压力可执行）。
\end{enumerate}
在该约束下，即便出现大量“迎合 MDQ/HACA-PACER”的微型模型，它们也只能在您定义的标准下竞争，而不会破坏体系的可移植性与治理闭环。
\end{remark}

\clearpage

%%%%%%%%%%%%%%%%%%%%%%%%%%%%%%%%%%%%%%%%%%%%%%%%%%%%%%%%%%%%
\section{白盒 AI 的程序回归：从端到端生成到“软提案 + 硬裁决”的谓词程序模式}
\label{sec:whitebox-program-regression}
%%%%%%%%%%%%%%%%%%%%%%%%%%%%%%%%%%%%%%%%%%%%%%%%%%%%%%%%%%%%

\begin{remark}[核心陈述（结构必然）]
在“LLM 只提案、MDQ 证书裁决、同伦类签名收敛”的结构下，白盒 AI 很可能出现一个明显的回归趋势：
从“端到端生成答案”回到“逻辑（谓词）驱动的程序模式”，但它不是回到早期手写规则的专家系统，
而是回到一种更现代、更可扩展的形态——\emph{模型负责提出候选程序，运行时负责验证与裁决}。
这不是口头约束，而是由可复验的裁判（证书谓词库 + 标准化器 + 参考执行器）与审计指标共同倒逼的结构结果（第\ref{sec:selection-pressure-llm-proposal-layer}节）。
\end{remark}

\subsection{最简结构：\texorpdfstring{$\mathrm{Comp}_m$}{Comp\_m} 产出候选“程序”，\texorpdfstring{$\mathrm{Eval}$}{Eval} 与
  \texorpdfstring{$V_s$}{V\_s} 裁决}
\label{subsec:program-minimal-structure}

沿用第\ref{subsec:comp-eval-contract}小节的记号：
\[
\mathrm{Comp}_m(x)\ \to\ \gamma\in\mathcal{D}^{\ast},
\qquad
\mathrm{Eval}(\gamma)\ \to\ s\in\mathcal{S}\sqcup\{\bot\}.
\]
其中 $\gamma$ 是一串 MDQ 调用（可理解为候选程序/候选证明链/候选编译产物），$\mathrm{Eval}$ 是确定性评估器。
对给定输入 $x$，最终“属于/不属于”由逻辑谓词输出，而非模型口头结论：
\[
\mathrm{Accept}(x\in s)\ \Longleftrightarrow\ V_s\bigl(C(x)\bigr)=1
\quad(\text{公理~\ref{ax:certificate-decision-priority}}).
\]

\begin{remark}[白盒的硬边界]
在该结构下，LLM 的自由度被显著收窄：它可以提不同的 $\gamma$，但只要进入裁决层，
就必须接受确定性复验与审计；任何不符合标准谓词边界的提案都会被拒绝并留痕（$\mathrm{FP}/\mathrm{FN}$，第\ref{subsec:fp-fn-cost}小节）。
\end{remark}

\subsection{语义本体压缩为可执行谓词：从同伦类到“组合谓词程序”}
\label{subsec:predicate-program}

\begin{definition}[MDQ 原子谓词]
对每个 MDQ 原子 $\delta=(\mathrm{span},\mathrm{type},V_\delta,R_\delta)\in\mathcal{D}$，定义其原子谓词为
\[
P_\delta:=V_\delta:\mathcal{X}_c\to\{0,1\}.
\]
\end{definition}

\begin{definition}[类谓词（同伦类谓词）]
对每个语义签名（同伦类）$s\in\mathcal{S}$，定义其类谓词为
\[
P_s:=V_s:\mathcal{X}_c\to\{0,1\}.
\]
\end{definition}

\begin{remark}[组合谓词的工程口径]
在工程层面，许多同伦类谓词可被理解为若干原子谓词的组合（此处用“近似”强调：具体组合口径由证书库定义）：
\[
P_s\ \approx\ \bigwedge_{i=1}^{k} P_{\delta_i}.
\]
组合方式不仅可以是合取，也可以包含析取、阈值化、约束组合与参数化谓词；关键是：它仍然是可复验的程序结构，而不是口头解释。
\end{remark}

\subsection{白盒解释的含义：从“能解释一句话”到“能回放一条证书链”}
\label{subsec:whitebox-as-replay}

\begin{definition}[证书链回放（最小口径）]
对白盒系统而言，一次接受/拒绝的最小解释证据是可回放证书链：
记录每一步调用了哪个 $\delta$（或哪个 $V_s$）、使用了哪个版本的验证器与标准化器、为何通过/失败、失败原因码是什么，
并可在同一 \code{verifier_version}/\code{policy_snapshot} 下复验出同结果（第\ref{subsec:replay-audit}小节）。
\end{definition}

\begin{remark}[更“硬”的白盒]
相比“用自然语言解释为什么”，证书链解释更硬：它可复验、可回放、可对比版本漂移（$H_{\mathrm{cert}}$，第\ref{subsec:certificate-implementation-hash}小节），
并且直接落在审计事件流上（第\ref{sec:mdq-granular-audit}节）。
\end{remark}

\subsection{LLM 的角色回归为“编译器前端/调度器”：优势在于更快找到可通过的 \texorpdfstring{$\gamma$}{gamma}}
\label{subsec:llm-as-frontend}

\begin{remark}[角色迁移]
在该体系中，LLM 的优势不再是“更会说”，而是“更快找到正确的 $\gamma$”：更小候选集、更少回退、更低成本、更低延迟。
系统会通过倒逼机制迫使提案器输出更贴合标准的 MDQ 调用：
不贴合就会在 $V_s(C(x))$ 处被拒绝并留痕，长期指标（$\mathrm{FP}/\mathrm{FN}$、回退率、延迟）会把它降权（第\ref{sec:selection-pressure-llm-proposal-layer}节）
  。
\end{remark}

\subsection{硬裁决并不等于只剩硬规则：稳定分工是“硬裁决 + 软提案”}
\label{subsec:hard-soft-separation}

\begin{remark}[工程平衡]
逻辑驱动并不等于“只剩硬规则”，而是形成稳定分工：
\begin{itemize}
  \item \textbf{硬的部分：}裁决与审计（$C$、$V_s$、$\mathrm{Eval}$）必须确定性、可复验、可治理；
  \item \textbf{软的部分：}提案与搜索（$\mathrm{Comp}_m$、PACER 策略、GRL 启发式）可以用模型，也可以迭代替换。
\end{itemize}
这使系统同时具备程序模式的稳定性与对任意风格输入的覆盖能力。
\end{remark}

\subsection{三值/部分逻辑：\texorpdfstring{$\{\mathrm{true},\mathrm{false},\bot\}$}{\{true,false,⊥\}} 更贴近工程事实}
\label{subsec:three-valued-logic}

\begin{definition}[部分裁决值域]
在工程实践中，严格的二值（true/false）往往不够表达“应当拒绝判定”的状态。
因此可将裁决值域扩展为三值：
\[
\{\mathrm{true},\ \mathrm{false},\ \bot\},
\]
其中 $\bot$ 表示“在当前证书基底与安全约束下不可判定/不应判定”。
\end{definition}

\begin{remark}[\texorpdfstring{$\bot$}{⊥} 的治理意义]
$\bot$ 避免模型硬猜，并把失败变成可治理信号：PACER 可将其解释为“需要补证书/需要升级 witness 档位/需要人工介入”，
并把原因码写入审计事件流，而不是把错误塞进 true/false（参见第\ref{subsec:basis-coverage-residual}小节的残余状态记账）。
\end{remark}

\subsection{与 OS+JVM 类比在此处更强：标准成为“逻辑运行时”，厂商竞争在提案器与优化器}
\label{subsec:osjvm-stronger}

\begin{remark}[生态形态]
当 Vol3 提供抽象规范（Spec），而本次讨论补齐运行时与一致性测试（Reference Runtime + TCK），
白盒 AI 的生态就会呈现出类似“语言/编译器/运行时/验证器/测试套件”的结构：
标准就是逻辑运行时（$C,V_s,\mathrm{Eval}$），厂商竞争在提案器与优化器（$\mathrm{Comp}_m$、PACER 策略），
并由同一裁判与 conformance suite 约束（第\ref{sec:os-jvm-paradigm}节）。
\end{remark}

\subsection{结构必然性：两条运行时公理立住后，“回归程序模式”不是趋势判断}
\label{subsec:structural-necessity}

\begin{remark}[两条公理作为落点]
若将以下两条运行时公理视为系统不变约束：
\begin{enumerate}
  \item \textbf{语义守恒（回投影不变性）：}公理~\ref{ax:projection-invariance}；
  \item \textbf{证书裁决优先：}公理~\ref{ax:certificate-decision-priority}；
\end{enumerate}
则系统中心就必然从“生成文本”迁移到“生成并通过验证的程序（证书调用序列）”：
因为任何可接受输出都必须能被 $\pi$ 回投影与 $V_s(C(x))$ 复验支撑，白盒性由可回放证书链给出。
因此“白盒回归程序模式”不再只是趋势判断，而是由结构约束推出的必然结果。
\end{remark}

\clearpage

%%%%%%%%%%%%%%%%%%%%%%%%%%%%%%%%%%%%%%%%%%%%%%%%%%%%%%%%%%%%
\section{对 Vol3 的增益评估：补齐运行时标准层与多实现一致裁决}
\label{sec:vol3-impact-evaluation}
%%%%%%%%%%%%%%%%%%%%%%%%%%%%%%%%%%%%%%%%%%%%%%%%%%%%%%%%%%%%

\begin{remark}[核心评价（对 Vol3 的直接意义）]
对 \code{docs/AppMathVol3/Pub_GFramework_AppMath_Vol3.tex}（Vol3）而言，本次讨论的意义非常直接：
它把 Vol3 已经出现的“语义规范场 / 语义函子场 / MDQ / HACA-PACER / 证书体系 / 治理与审计”主线，
补上一个\emph{可执行、可移植、可度量}的运行时标准层。换句话说：在不改变 Vol3 的理论骨架前提下，
把“语义地板”以工程方式抬上去，并给出一套能在多模型、多实现之间保持一致的裁决机制。
\end{remark}

本章按“Vol3 已有内容 $\rightarrow$ 本次讨论新增关键增益 $\rightarrow$ 可对接到 Vol3 的哪一块”来归纳与评价。

\subsection{增益 1：为 PACER 固化可执行最小合同——压缩/展开/回投影不变}
\label{subsec:vol3-gain-1-roundtrip}

Vol3 的 PACER 叙事以“压缩/展开推理”为主线。本次讨论将其固化为一个硬条件（可作为系统公理与一致性测试）：
\[
s=\pi(x),\qquad x_\sigma = G(s,\sigma),\qquad \pi(x_\sigma)=s\ \ (\forall \sigma).
\]
其中 $\pi$ 是语义投影/同伦类裁决（由 MDQ 证书体系确定），$G$ 是按风格 $\sigma$ 展开。
这正是公理~\ref{ax:projection-invariance} 的运行时化表达：把“不失真”从写作主张升级为可测规范，
并天然支持“多风格冗余展开”作为不失真证据（第\ref{subsec:lossless-expansion}小节）。

\subsection{增益 2：把 MDQ 具体化为证书化可调用单元，从而支持 MDQ 粒度审计}
\label{subsec:vol3-gain-2-mdq-cert}

Vol3 已引入 MDQ（differential dynamical quanta）用于路径积分与引擎结构。本次讨论的关键补强是把每个 MDQ 具体化为可外包证书：
\[
\delta = (\mathrm{span},\mathrm{type},V_\delta,R_\delta).
\]
其中 $V_\delta$ 是确定性验证器（正则/自动机/小解析器等的受限子集），$R_\delta$ 是标准化重写；
$\mathrm{span}/\mathrm{type}$ 把责任定位到最小片段与类型（第\ref{subsec:mdq-basis-delta}小节）。
因此审计可下沉到“哪个 $\delta$ 出错、哪条证书版本漂移、哪个 verifier 失败”，而不是停留在整段文本层面的模糊追责（第\ref{sec:mdq-granular-audit}节）。

\subsection{增益 3：LLM 只提案，证书来裁决——把语义正确性从概率叙事变成可复验事实}
\label{subsec:vol3-gain-3-llm-propose}

对 Vol3 的工程意义在于：把 LLM 放在正确位置——不是裁判，而是编译器/提案器。
设提案为 $(\hat{s},\hat{y})$，确定性复验为
\[
y^\ast = V_{\hat{s}}(C(x)).
\]
则两类核心错误可机械定义并留痕：
误收（$\hat{y}=1 \wedge y^\ast=0$）与漏收（$\hat{y}=0 \wedge y^\ast=1$），对应第\ref{subsec:fp-fn-cost}小节的 $\mathrm{FP}/\mathrm{FN}$
  记账。
这使 Vol3 的“证书 + 解释 + 治理”不再依赖模型自觉，而依赖可复验谓词裁决（公理~\ref{ax:certificate-decision-priority}），
从而直接补强“版本对比/审计证据/回滚决策”等管线的可信度（第\ref{subsec:replay-audit}小节）。

\subsection{增益 4：“多次扫描后语义明确”的严格含义——对同伦类（语义签名）收敛}
\label{subsec:vol3-gain-4-convergence}

本次讨论将“收敛”明确为：候选同伦类集合单调收缩并达到固定点（第\ref{subsec:candidate-set-scan-operator}小节）。
工程上这意味着：
\begin{itemize}
  \item 收敛到单点：语义明确（标准语义锚点成立）；
  \item 固定点 $>1$ 或空集：不是玄学，而是“证书基底覆盖不足/同伦边界定义不足”的可定位缺口（第\ref{subsec:basis-coverage-residual}小节）。
\end{itemize}
因此 Vol3 中“语义规范场可修复表达上同调障碍”的叙事获得工程化判据：障碍可被记账为“固定点不存在/不稳定”，从而可被 PACER 的修复策略系统化处理。

\subsection{增益 5：同伦类哈希与证书链哈希——把证书治理升级为防伪锚点（并带出防洗稿能力）}
\label{subsec:vol3-gain-5-hash}

Vol3 已讨论证书/合规/治理。本次讨论补出的关键资产是稳定语义指纹与可回放实现指纹：
\[
H_{\mathrm{sem}}(x)=\operatorname{Hash}\bigl(\mathrm{ClassID}(\pi_0(x))\bigr),
\qquad
H_{\mathrm{cert}}(x)=\operatorname{Hash}\bigl(\mathrm{ClassID},\mathrm{PredicateID},\mathrm{WitnessHash},
  \mathrm{VerifierVer},\mathrm{CanonVer}\bigr).
\]
$H_{\mathrm{sem}}$ 解决“语义层稳定索引/防伪/聚类/去重”；$H_{\mathrm{cert}}$ 解决“版本漂移解释、证书篡改不可否认”（第\ref{sec:homotopy-fingerprint-hash}节）。
在该口径下，“洗稿式改写”在数学上就是同伦类内的风格扰动；您既能在同一 $\pi(x)$ 内展开多风格表达，也能用 $H_{\mathrm{sem}}$ 与局部指纹把它识别出来并审计（第\ref{subsec:layer-local-fingerprint}小节）。

\subsection{增益 6：OS+JVM 类比不是修辞——把 Vol3 推向“标准—多实现—一致性测试”的生态形态}
\label{subsec:vol3-gain-6-osjvm}

Vol3 若要走向“多实现可互换”的生态形态，需要一个类似 \textsf{JVM TCK} 的一致性测试框架。
本次讨论把该生态层写成形式合同：对任意实现 $m$，令其生成 MDQ 调用序列 $\mathrm{Comp}_m(x)$，由参考执行器裁决 $\mathrm{Eval}$，
则必须满足命题~\ref{prop:portability-contract} 所刻画的可移植性合同（可验证时回到同一 $\pi_0(x)$）：
\[
\mathrm{Eval}\bigl(\mathrm{Comp}_m(x)\bigr)=\pi_0(x).
\]
配合 “Spec + Reference Verifier + Conformance Suite”（第\ref{subsec:spec-verifier-tck}小节），该表达显著增强 Vol3 对工程读者（系统/语言/平台/合规）
  的可理解度与可采纳度。

\subsection{增益 7：对 GRL 路径积分“发现新语义/意识空间”的意义——把探索空间离散化为可证书化轨线}
\label{subsec:vol3-gain-7-grl}

Vol3 的 GRL path integral 已把学习写成路径积分。本次讨论补强的是：轨线每一步是 MDQ 证书化步进 $\delta_t$，
因此“探索新语义”不再是不可审计生成，而是可重放的路径搜索：
\[
s_{t+1}=s_t \oplus \delta_t,\qquad \delta_t\in\mathcal{D}.
\]
再配合作用量（歧义/成本/风险/新颖性）对路径加权（第\ref{subsec:grl-path-integral}小节），即可把“稳定出现的新终点”固化为新同伦类与新证书。
于是“意识空间”（建模意义）具备：状态（语义签名/分布）+ 算子（MDQ/变换）+ 可审计轨线（证书链）（第\ref{subsec:consciousness-space-modeling}小节）。

\subsection{对接点：建议落到 Vol3 的三处文本块（点名对齐）}
\label{subsec:vol3-docking-points}

\begin{remark}[三处集中对接点]
从文本对接角度，本次讨论最有价值的落点可集中在三处（不要求立即改稿，但便于定位意义）：
\begin{itemize}
  \item \textbf{语义规范场与 triple-layer connection 的联络/曲率叙事处：}
  本次讨论提供“离散协变微分 = MDQ 证书调用”的工程释义（第\ref{subsec:semantic-gauge-field-connection}小节）。
  \item \textbf{Law-Space Certificates \& MDQ-Based Evaluation 与 System Architectures/Safety/Governance：}
  本次讨论补上“误收/漏收留痕 + 哈希指纹 + conformance suite + 多厂商一致性”的可执行框架（第\ref{sec:mdq-granular-audit}、\ref{sec:homotopy-fingerprint-hash}、\ref{sec:os-jvm-paradigm}节）。
  \item \textbf{PACER 的压缩/展开主线：}
  本次讨论给出回投影不变性作为硬验收标准，使“通俗展开、多风格展开、语义不失真”从叙事升级为可验证条件（公理~\ref{ax:projection-invariance}）。
\end{itemize}
\end{remark}

\subsection{综合评价（摘要）}
\label{subsec:vol3-summary}

\begin{remark}[一句话总结]
本次讨论把 Vol3 的核心思想从“描述一个可治理的语义体系”，推进到“给出一套可被多方实现、用一致性测试约束、并通过哈希与证书链实现可审计与防伪的运行时标准”：
学术上增强可检验性，工程上增强可落地性与可移植性，生态上给出清晰的标准化路径。
\end{remark}

\clearpage

%%%%%%%%%%%%%%%%%%%%%%%%%%%%%%%%%%%%%%%%%%%%%%%%%%%%%%%%%%%%
\section{分层结论：Vol3 作为 Spec，本次讨论作为 Reference Runtime + Conformance Suite}
\label{sec:spec-vs-runtime-layering}
%%%%%%%%%%%%%%%%%%%%%%%%%%%%%%%%%%%%%%%%%%%%%%%%%%%%%%%%%%%%

\begin{remark}[核心划分（写作与落地更干净）]
将体系分为“应用数学规范层（Vol3）”与“工程运行时标准层（本次讨论）”，对后续写作与落地更有利：
\begin{itemize}
  \item \textbf{Vol3 的主任务：}给出“语义规范场 / 语义函子 / MDQ / HACA-PACER / GRL 路径积分”的抽象对象、等价关系、算子结构与推导链条，
  强调定义严谨性、可证明性、可复用的抽象接口（应用数学层面的规范/定理系统）。
  \item \textbf{本次讨论的主任务：}把抽象对象落成一套“能跑、能审计、能版本治理、能多厂商实现”的工程机制，
  强调可执行合同、裁决闭环、性能与容错指标、实现可移植性（工程层面的运行时/参考实现/一致性测试体系）。
\end{itemize}
\end{remark}

\subsection{工程化对应：Spec 与 Reference Runtime 的分工}
\label{subsec:spec-runtime-roles}

\begin{remark}[对应关系]
\begin{itemize}
  \item Vol3 更像 \textbf{规范（Spec）}：规定语义对象是什么、等价是什么、算子如何作用、收敛/修复应满足什么性质；
  \item 本次讨论更像 \textbf{参考运行时（Reference Runtime）+ 一致性测试（Conformance Suite）}：
  规定如何把 $\pi$、$V_s$、MDQ、审计、哈希、回投影不变性等条件做成可执行裁判与可回放证据链。
\end{itemize}
\end{remark}

\subsection{OS+JVM 类比在分层处的落点}
\label{subsec:osjvm-layering}

\begin{remark}[类比成立的原因]
在第\ref{sec:os-jvm-paradigm}节的 OS+JVM 类比下：
\begin{itemize}
  \item Vol3 负责把“语义中间层”的语义规范写清楚（类似字节码语义规范）；
  \item 本次讨论负责把“验证器 + 审计 + 测试套件 + 多实现一致”补齐（类似 bytecode verifier + TCK + 多实现生态）。
\end{itemize}
因此该类比不是修辞，而是对“标准—多实现—一致性测试”生态结构的分层说明。
\end{remark}

\subsection{运行时要件清单：本次讨论为 Vol3 补齐的不可缺部件}
\label{subsec:runtime-requirements}

\begin{remark}[七类运行时要件（与前文一一对应）]
纯应用数学写法通常不会展开运行时细节；本次讨论为工程实现层补齐了至少七类不可缺部件：
\begin{enumerate}
  \item \textbf{裁决机制：}LLM 只提案，证书谓词 $V_s$ 裁决（公理~\ref{ax:certificate-decision-priority}）；
  \item \textbf{MDQ 证书化：}把 MDQ 落成可调用、可版本化、可回放的最小单元（第\ref{subsec:mdq-basis-delta}小节）；
  \item \textbf{收敛判据：}多次扫描等价于对同伦类签名（及其哈希）达到固定点（第\ref{sec:multi-scan-convergence}节）；
  \item \textbf{审计体系：}$\mathrm{FP}/\mathrm{FN}$ 机械可判定、可留痕、可统计（第\ref{subsec:fp-fn-cost}小节）；
  \item \textbf{防伪锚点：}$H_{\mathrm{sem}}$（语义指纹）与 $H_{\mathrm{cert}}$（证书链指纹）（第\ref{sec:homotopy-fingerprint-hash}节）；
  \item \textbf{一致性测试：}把回投影不变性做成工程验收标准（公理~\ref{ax:projection-invariance} 与第\ref{subsec:roundtrip-regression}小节）；
  \item \textbf{多实现可移植：}不同模型/厂商可优化提案策略，但必须通过同一裁判与测试套件（命题~\ref{prop:portability-contract} 与第\ref{subsec:spec-verifier-tck}
    小节）。
\end{enumerate}
\end{remark}

\subsection{写作与发布建议：理论—工程的互相加固闭环}
\label{subsec:publishing-suggestion}

\begin{remark}[建议]
若保持该分层，后续写作会更干净：Vol3 保持抽象与严谨；工程部分可作为独立的
\code{Implementation Spec / Reference Runtime} 文档或伴随仓库发布，并反过来用一致性测试与审计数据去验证 Vol3 的关键公理（例如回投影不变性与收敛性），形成理论—工程互相加固的闭环。
\end{remark}

\clearpage

\appendix

%%%%%%%%%%%%%%%%%%%%%%%%%%%%%%%%%%%%%%%%%%%%%%%%%%%%%%%%%%%%
\section{附录A：可对外调用的证书接口（MDQ Certificate API）}
\label{app:mdq-certificate-api}
%%%%%%%%%%%%%%%%%%%%%%%%%%%%%%%%%%%%%%%%%%%%%%%%%%%%%%%%%%%%

\begin{remark}[附录目标]
本附录给出一份\emph{可对外调用}且\emph{可复验}的证书接口合同，用于把本文的核心结构落到工程实现：
外部系统只需调用该接口即可获得“证书化的判别器 + 标准化重写 + 指纹与审计字段”，并可用参考验证器复验其有效性。
\end{remark}

\subsection{接口定位：对外只提供“证书裁决”，不提供口头解释}
\label{subsec:api-positioning}

\begin{definition}[信任模型（最小口径）]
对外接口的信任模型为：
\begin{enumerate}
  \item 客户端\emph{不需要}信任 LLM 的自然语言解释；
  \item 客户端\emph{只需要}信任确定性裁决链：标准化器 $C$、证书谓词 $V$（或 $V_s,V_\delta$）与参考执行/校验器；
  \item 任何返回为 \code{ok:true} 的结果，都必须携带足以复验的版本字段与摘要字段（第\ref{sec:homotopy-fingerprint-hash}节）。
\end{enumerate}
\end{definition}

\begin{remark}[对齐本文公理]
对外接口的“接受”语义由公理~\ref{ax:certificate-decision-priority} 给出：
\[
\mathrm{Accept}(x\in s)\Longleftrightarrow V_s(C(x))=1,
\]
并在通过时返回可复验证书；不通过则返回 $\bot$（实现上为 \code{ok:false}）。
\end{remark}

\subsection{传输与版本：HTTP+JSON、版本化与幂等}
\label{subsec:api-transport-versioning}

\begin{definition}[传输约定（建议）]
外部调用采用 \textbf{HTTP(S)+JSON}：
\begin{itemize}
  \item \code{Content-Type: application/json; charset=utf-8}；
  \item JSON 语义以 RFC~8259 为准\cite{rfc8259}；
  \item 所有哈希摘要以 \textsf{SHA-256} 为推荐默认\cite{fips180-4}（以 \code{sha256:...} 表示）。
\end{itemize}
\end{definition}

\begin{definition}[版本字段（必须）]
每次返回必须携带以下最小版本字段，以支持回放与漂移解释：
\begin{itemize}
  \item \code{api_version}：接口版本（如 \code{"mdq-cert-api.v1"}）；
  \item \code{verifier_version}：参考验证器版本；
  \item \code{canon_version}：标准化/规范化规则版本（含 $C$ 与 canonicalization）；
  \item \code{policy_snapshot_id}：白名单与复杂度阈值策略快照标识（第\ref{subsec:audit-objects}小节）。
\end{itemize}
\end{definition}

\begin{definition}[幂等与关联字段（建议）]
建议客户端为每次调用提供：
\begin{itemize}
  \item \code{request_id}：请求幂等键（用于重试与去重）；
  \item \code{trace_id}：跨系统链路追踪标识（可选）；
  \item \code{input_digest}：输入摘要（可选；若未提供，服务端可返回其计算值）。
\end{itemize}
\end{definition}

\subsection{核心数据结构：\texorpdfstring{$\mathsf{SemSig}$}{SemSig}、证书与证书链}
\label{subsec:api-core-objects}

\begin{definition}[\texorpdfstring{$\mathsf{SemSig}$}{SemSig}（对外签名对象）]
对外的语义签名对象建议采用：
\begin{verbatim}
{
  "class_id": "mdq.uuid",
  "params": { },
  "h_sem":  "sha256:...",
  "h_cert": "sha256:..."
}
\end{verbatim}
其中 \code{class_id} 对应 $\mathrm{ClassID}$，\code{h_sem} 与 \code{h_cert} 分别对应 $H_{\mathrm{sem}}$ 与 $H_{\mathrm{cert}}$
（第\ref{sec:homotopy-fingerprint-hash}节）。
\end{definition}

\begin{definition}[证书对象（最小可复验口径）]
对外证书对象建议封装为三元组（Spec/Witness/Verifier）在实现层的投影：
\begin{verbatim}
{
  "cert_id": "mdq.uuid.v1",
  "mdq_type": "UUID",
  "spec_hash": "sha256:...",
  "witness": {
    "kind": "regex",
    "payload": "^[0-9a-fA-F]{8}-...$"
  },
  "witness_hash": "sha256:...",
  "verifier_version": "verifier.v3",
  "canon_version": "canon.v2",
  "test_digest": "sha256:..."
}
\end{verbatim}
其中：
\begin{itemize}
  \item \code{spec_hash} 固化语义边界（Spec）的摘要；
  \item \code{witness} 为判别器本体（正则/自动机/小解析器三档之一，见第\ref{subsec:mdq-equivalence}小节）；
  \item \code{witness_hash} 固化判别器实现；
  \item \code{test_digest} 固化回归/对抗测试集的摘要（用于审计回放）。
\end{itemize}
\end{definition}

\begin{definition}[证书链（链式输出）]
若一次裁决涉及多个 MDQ 单元，则返回证书链：
\begin{verbatim}
"cert_chain": [ {cert_1}, {cert_2}, ... ]
\end{verbatim}
并要求链中每个证书都可独立复验；链整体用于定位失败点、做差分审计与回放（第\ref{sec:mdq-granular-audit}节）。
\end{definition}

\subsection{接口 1：签发/裁决（\texorpdfstring{\texttt{POST /v1/certify}}{POST /v1/certify}）}
\label{subsec:api-certify}

\begin{definition}[请求：\code{CertifyRequest}]
请求体（示意）：
\begin{verbatim}
{
  "api_version": "mdq-cert-api.v1",
  "request_id": "uuid-...",
  "input": { "expr": "...", "lang": "zh-CN" },
  "policy_snapshot_id": "policy.safe-re2.v1",
  "options": { "return_examples": false, "return_audit": true }
}
\end{verbatim}
\end{definition}

\begin{definition}[响应：\code{CertifyResponse}]
响应体满足二值输出合同（\code{ok:true}+证书 / \code{ok:false}）：
\begin{verbatim}
// ok:true
{
  "ok": true,
  "api_version": "mdq-cert-api.v1",
  "request_id": "uuid-...",
  "semsig": { ... },
  "certificate": { ... },
  "cert_chain": [ ... ],
  "verifier_version": "verifier.v3",
  "canon_version": "canon.v2",
  "policy_snapshot_id": "policy.safe-re2.v1"
}

// ok:false
{
  "ok": false,
  "api_version": "mdq-cert-api.v1",
  "request_id": "uuid-...",
  "reason_code": "WHITELIST_VIOLATION",
  "policy_snapshot_id": "policy.safe-re2.v1"
}
\end{verbatim}
\end{definition}

\begin{remark}[失败原因码（建议枚举）]
建议将失败原因码限定为稳定枚举，便于审计聚合与回归定位，例如：
\code{JSON_INVALID}、\code{REGEX_COMPILE_FAIL}、\code{WHITELIST_VIOLATION}、
\code{MISSING_ANCHORS}、\code{EXAMPLE_CONTRADICTION}、
\code{COMPLEXITY_LIMIT}、\code{METAMORPHIC_REGRESSION_FAIL}、
\code{POLICY_REJECT}、\code{NON_REGULAR_SEMANTICS} 等。
\end{remark}

\subsection{接口 2：复验（\texorpdfstring{\texttt{POST /v1/verify}}{POST /v1/verify}）}
\label{subsec:api-verify}

\begin{definition}[请求：\code{VerifyRequest}]
\begin{verbatim}
{
  "api_version": "mdq-cert-api.v1",
  "request_id": "uuid-...",
  "input": { "expr": "...", "lang": "zh-CN" },
  "certificate": { ... },
  "policy_snapshot_id": "policy.safe-re2.v1"
}
\end{verbatim}
\end{definition}

\begin{definition}[响应：\code{VerifyResponse}]
\begin{verbatim}
{
  "ok": true,
  "request_id": "uuid-...",
  "verify_result": true,
  "reason_code": "OK",
  "verifier_version": "verifier.v3",
  "policy_snapshot_id": "policy.safe-re2.v1"
}
\end{verbatim}
\end{definition}

\begin{remark}[验证语义]
复验接口必须是确定性的：在固定 \code{verifier_version} 与 \code{policy_snapshot_id} 下，
同一 \code{certificate} 与同一输入应产生一致的 \code{verify_result}；
这对应第\ref{subsec:replay-audit}小节的“可回放一致性”要求。
\end{remark}

\subsection{最小对外保证：可移植、可回放、可对比}
\label{subsec:api-min-guarantees}

\begin{proposition}[对外可复验性（合同）]
若 \code{CertifyResponse.ok=true}，则客户端可以在相同的 \code{verifier_version}、\code{canon_version} 与 \code{policy_snapshot_id} 下，
通过 \code{POST /v1/verify} 复验通过（\code{verify_result=true}），并且该结果与证书链指纹 $H_{\mathrm{cert}}$ 一致。
\end{proposition}

\begin{remark}[对齐“标准先行、多实现”】【工程解释】]
该附录接口为“标准先行、多实现”（第\ref{sec:os-jvm-paradigm}节）提供一个最小可执行落点：
上层可以有多种提案实现（不同 LLM/不同策略），但一旦对外输出为 \code{ok:true}，
其证据必须能被同一参考验证器复验成立，从而实现可移植与可审计。
\end{remark}

\begin{thebibliography}{99}
\addcontentsline{toc}{section}{\refname}

\bibitem{rfc8259}
T.~Bray.
\emph{The JavaScript Object Notation (JSON) Data Interchange Format}.
RFC~8259, IETF, 2017.
\url{https://www.rfc-editor.org/rfc/rfc8259}

\bibitem{cox-re2}
R.~Cox.
\emph{RE2: a fast, safe, thread-friendly regular expression engine}.
Google.
\url{https://github.com/google/re2}

\bibitem{cox-regex}
R.~Cox.
\emph{Regular Expression Matching Can Be Simple And Fast}.
2007.
\url{https://swtch.com/~rsc/regexp/regexp1.html}

\bibitem{chen-metamorphic}
T.~Y.~Chen, S.~C.~Cheung, S.~M.~Yiu.
\emph{Metamorphic Testing: a New Approach for Generating Next Test Cases}.
Technical Report HKUST-CS98-01, 1998.

\bibitem{hopcroft-ullman}
J.~E.~Hopcroft, R.~Motwani, J.~D.~Ullman.
\emph{Introduction to Automata Theory, Languages, and Computation}.
3rd~ed., Addison-Wesley, 2006.

\bibitem{fips180-4}
NIST.
\emph{FIPS PUB 180-4: Secure Hash Standard (SHS)}.
2015.
\url{https://csrc.nist.gov/publications/detail/fips/180/4/final}

\bibitem{jvms}
T.~Lindholm, F.~Yellin, G.~Bracha, A.~Buckley.
\emph{The Java Virtual Machine Specification, Java SE 8 Edition}.
Addison-Wesley, 2015.

\bibitem{java-tck}
Oracle.
\emph{Java SE Technology Compatibility Kit (TCK)}.
\url{https://www.oracle.com/java/technologies/javase-tck.html}

\bibitem{GaoZhengAppVol3}
Gao, Z. (2025).
\newblock GaoZheng G-Framework and GaoZheng G-Algebra: Applied Mathematics Volume III – HACA, PACER, and
  Certificate-Based AI.
\newblock In \emph{GaoZheng G-Framework and GaoZheng G-Algebra: Applied Mathematics Volume III – HACA, PACER, and
  Certificate-Based AI (v1.0)}.
\newblock Zenodo.
\newblock \url{https://doi.org/10.5281/zenodo.17762295}.


\bibitem{gaozheng-analytic-ai}
GaoZheng.
\emph{基于主纤维丛版广义非交换李代数的解析解AI系统设计}.
内部文档，2025。

\end{thebibliography}

\end{CJK*}
\end{document}
